% ============================================================
%  Measuring What We Promise: A Cross-Institutional Analysis
%  of Learning Outcome Quality in Irish Higher Education
%
%  Draft v0.1 — April 2026
%  Authors: Keating, Fortune, Gottlöber
% ============================================================

\documentclass[12pt, a4paper]{article}

\usepackage[margin=2.5cm]{geometry}
\usepackage[utf8]{inputenc}
\usepackage[T1]{fontenc}
\usepackage{lmodern}
\usepackage{microtype}
\usepackage{setspace}
\usepackage{booktabs}
\usepackage{tabularx}
\usepackage{array}
\usepackage{multirow}
\usepackage{graphicx}
\usepackage{amsmath}
\usepackage{hyperref}
\usepackage{natbib}
\usepackage{xcolor}
\usepackage{enumitem}
\usepackage{caption}
\usepackage{float}
\usepackage{csquotes}

\hypersetup{
  colorlinks=true,
  linkcolor=black,
  citecolor=black,
  urlcolor=blue
}

\setstretch{1.5}
\setlength{\parskip}{6pt}
\setlength{\parindent}{0pt}

% ── Draft annotation colour ─────────────────────────────────
\newcommand{\todo}[1]{\textcolor{red}{\textbf{[TODO: #1]}}}
\newcommand{\note}[1]{\textcolor{blue}{\textit{[Note: #1]}}}

% ────────────────────────────────────────────────────────────
\begin{document}

\begin{titlepage}
  \centering
  \vspace*{2cm}

  {\LARGE\bfseries Measuring What We Promise:\\[0.4em]
  A Cross-Institutional Analysis of Learning Outcome\\[0.2em]
  Quality in Irish Higher Education\par}

  \vspace{1.5cm}

  {\large
    John G. Keating$^{1}$,
    Niamh Fortune$^{2}$,
    Susan Gottl{\"o}ber$^{3}$,
    M{\'a}ire Buckley$^{4}$,
    Tim Thompson$^{5}$,
    Lisa O'Regan$^{6}$
  }

  \vspace{0.8cm}

  {\normalsize
    $^{1}$Associate Dean for Teaching and Learning, Faculty of Science and Engineering\\
    $^{2}$Associate Dean for Teaching and Learning, Faculty of Social Sciences\\
    $^{3}$Associate Dean for Teaching and Learning, Faculty of Arts and Humanities\\
    $^{4}$Experiential Learning Officer, Student Skills and Success\\
    $^{5}$Vice President for Students and Learning\\
    $^{6}$Head, Centre for Teaching and Learning\\[0.4em]
    Maynooth University
  }

  \vspace{0.8cm}

  {\normalsize
    Correspondence: \texttt{john.keating@mu.ie}
  }

  \vspace{1.5cm}

  {\normalsize \textit{Draft submitted for review --- \today}}

\end{titlepage}

% ── Abstract page ────────────────────────────────────────────
\newpage
\begin{abstract}
  \noindent
  Learning outcomes (LOs) are the primary accountability instrument in Irish higher
  education, underpinning programme approval, Quality and Qualifications Ireland (QQI)
  compliance, and the National Framework of Qualifications (NFQ). Yet the quality of
  LO language --- specifically the observability and cognitive demand of the action verbs
  employed --- has received little systematic empirical attention at scale. This paper
  reports on a corpus-scale analysis of \textbf{233,005} module-level LOs drawn from
  \textbf{sixteen} Irish higher education institutions (HEIs), spanning the public university,
  technological university (TU), Institute of Technology (IT), College of Education, and private/specialised sectors. Using a rule-based
  natural language processing (NLP) pipeline aligned to a six-dimension evaluation
  rubric (D1--D6), we score each LO on verb observability (D1), behavioural singularity
  (D2), and student-centredness (D3). We report a consistent and statistically meaningful
  gap between the public university and TU clusters on the D1 (verb observability)
  dimension: public university mean D1 = 1.572 versus TU mean D1 = 1.775
  ($\Delta$ = 0.203), with public universities showing substantially higher rates
  of unobservable Tier~0 verbs (8.3\% versus 3.4\%). We extend the analysis to
  constructive alignment, cross-tabulating verb tier against assessment modality and,
  for three institutions with Akari curriculum management data, against LO-to-component
  mappings, identifying 64 modules in which high-demand LOs are assessed exclusively
  by formal examination. We discuss implications for institutional quality assurance
  processes, module approval workflows, and national policy. Preliminary findings from
  Universal Design for Learning (UDL) proxy analysis and HEA Data and Digital Skills
  module tagging are also presented.
\end{abstract}

\vspace{0.5cm}
\noindent\textbf{Keywords:} learning outcomes; higher education quality; natural language
processing; constructive alignment; verb taxonomy; Irish higher education; QQI; NFQ;
Universal Design for Learning
\newpage

% ════════════════════════════════════════════════════════════
\section{Introduction}
% ════════════════════════════════════════════════════════════

The learning outcome has become the dominant currency of higher education quality
assurance in Europe and Ireland. Across national qualifications frameworks, programme
approval processes, institutional review, and external examiner reports, the question
asked --- implicitly or explicitly --- is the same: what can a graduate \emph{do}?
The shift from input-based to outcome-based education is now decades old
\citep{adam2004using, kennedy2007writing}, yet a persistent tension remains between
the aspiration embedded in that question and the language used to answer it.

If a module learning outcome states that a student will \textit{understand} the
principles of thermodynamics, or will \textit{appreciate} the historical context of
modernism, what exactly has been promised? The verb carries the epistemic weight of
the outcome statement, yet verbs of this kind --- vague, unobservable, unmeasurable
--- remain prevalent in published module descriptors across the sector
\citep{hussey2002trouble, moon2002module}. They satisfy the formal requirement of
outcome-based language while resisting the accountability that language was designed
to enable.

This paper addresses that gap directly. We report on a systematic, corpus-scale
analysis of module-level learning outcomes drawn from fourteen Irish HEIs: six public
universities, three technological universities, two Institutes of Technology, one College
of Education, and two private/specialised colleges.
Our corpus of 229,037 LOs was assembled through
systematic retrieval of publicly accessible institutional module catalogues, processed through a rule-based NLP
pipeline, and evaluated against a six-dimension rubric developed from the literature
on LO quality, Bloom's revised taxonomy \citep{anderson2001taxonomy}, and the Irish
National Framework of Qualifications \citep{qqi2020nfq}.

The research makes three substantive contributions. First, it provides the first
cross-institutional, corpus-scale empirical characterisation of LO language quality
in Irish higher education. Second, it identifies a consistent and reproducible gap
between the university and technological university clusters on the dimension of
verb observability. Third, it extends the analysis to constructive alignment
\citep{biggs2011teaching}, demonstrating that verb quality alone is insufficient:
high-demand LOs assessed exclusively by traditional formal examination represent
a distinct and actionable category of misalignment.

The paper is structured as follows. Section~\ref{sec:background} reviews the relevant
literature. Section~\ref{sec:methods} describes corpus construction, the evaluation
rubric, and the scoring pipeline. Section~\ref{sec:results} presents findings across
all nine institutions with particular focus on the university--TU comparison and
constructive alignment. Section~\ref{sec:discussion} interprets findings in the context
of Irish HE policy and institutional practice. Section~\ref{sec:conclusion} offers
recommendations.

% ════════════════════════════════════════════════════════════
\section{Background and Related Literature}
\label{sec:background}
% ════════════════════════════════════════════════════════════

\subsection{Learning Outcomes in Higher Education Policy}

The outcomes-based approach to curriculum design entered Irish higher education
formally through the Bologna Process and the establishment of the National Framework
of Qualifications in 2003 \citep{nfq2003}. QQI's \textit{Policies, Criteria and
Guidelines for the Validation of Programmes of Education and Training}
\citep{qqi2021policies} require that all accredited programmes express intended
learning outcomes that are specific, assessable, and aligned to NFQ award standards.
The National Forum for the Enhancement of Teaching and Learning has further endorsed
outcome-based design as a cornerstone of curriculum quality
\citep{nationalforum2016mapping}.

Despite this policy consensus, empirical evidence on the \emph{quality} of LO
language in practice is sparse. Institutional review bodies assess whether outcomes
are present; they do not routinely measure their linguistic or cognitive quality.
\citet{hussey2002trouble} identified early that outcome statements frequently conflate
multiple behaviours, use unobservable language, or are pitched at levels inconsistent
with their NFQ designation. \citet{moon2002module} similarly noted the gap between
outcome-as-written and outcome-as-assessable. More recently, \citet{donnelly2019outcomes}
observed that Irish HEIs have made significant structural progress in outcome adoption
but that the quality of individual statements remains uneven.

\subsection{The Policy Framework for Learning Outcomes: Three Levels}
\label{sec:policyframework}

\note{\textbf{[PLACEHOLDER: to be developed with Lisa O'Regan, Centre for Teaching
and Learning, MU. See accompanying policy reference document.]}}

LO requirements for Irish HEIs operate at three intersecting levels. At the
\textit{European} level, the Standards and Guidelines for Quality Assurance in the
European Higher Education Area (ESG 2015) \citep{enqa2015esg} require under
Standard~1.2 that programmes are \enquote{designed so that they meet the objectives
set for them, including the intended learning outcomes} and that \enquote{the
qualification resulting from a programme should be clearly specified and communicated,
and refer to the correct level of the national qualifications framework.} The Dublin
Descriptors, which underpin the Bologna Process qualification cycles, specify the
knowledge, skills, judgement, communication, and learning-to-learn capabilities
expected at each cycle, and implicitly require LO language that reaches the cognitive
demand appropriate to the cycle (Bachelor, Master, Doctorate) \citep{joint2004dublin}.

At the \textit{national} level, QQI's \textit{Policies, Criteria and Guidelines for
the Validation of Programmes of Education and Training} \citep{qqi2021policies} require
that learning outcomes are \enquote{specific and sufficiently detailed to communicate
the award-holder's knowledge, skill and competence} and that they are \enquote{appropriate
to the specified NFQ level and award type.} The NFQ Grid of Level Indicators defines
expected knowledge, skills, and competence at each level; at Level~8 (Honours
Bachelor) this includes \enquote{mastery of a complex and specialised area of skills
and tools} and the capacity to \enquote{conduct guided research} --- language that
maps to observable, high-demand verbs in the D1 Tier 2--3 range. QQI further requires
that \enquote{assessment procedures should validly and reliably ensure that learners
who are recommended for awards have attained the provider's minimum intended programme
learning outcomes} --- a direct statement of the constructive alignment requirement
\citep{qqi2021policies}. The National Forum for the Enhancement of Teaching and
Learning has further endorsed outcome-based design and multi-modal assessment (OF,
FOR, and AS learning) as cornerstones of curriculum quality
\citep{nationalforum2016mapping}.

At the \textit{institutional} level, each HEI is required under QQI's quality
assurance framework to maintain internal procedures for programme and module design
and approval that operationalise these requirements. The gap between this procedural
requirement and the documented outcome quality revealed by corpus analysis is a
central concern of this study.

Critically, the policy framework does not mandate Bloom's taxonomy by name. It does,
however, require precisely what Bloom formalised: outcomes that are specific,
observable, level-appropriate, and assessable. Our D1--D6 rubric operationalises
the policy criteria without depending on any single pedagogical framework: D1
(verb observability) and D2 (behavioural singularity) address the \enquote{specific
and sufficiently detailed} requirement; D3 (student-centredness) addresses the
student-centred LO format implied by the Bologna student-learning paradigm; D6
(NFQ level calibration) addresses the \enquote{appropriate to the specified NFQ
level} requirement. The rubric is, in this sense, a policy-compliance instrument
as much as a pedagogical quality instrument.

\todo{Lisa O'Regan to develop: (a) institutional policy context at MU (L\&T
Strategy, Programme Board procedures); (b) national QA framework developments
post-2021; (c) comparative European context — how does Irish policy implementation
compare with EHEA peers? (d) implications of the shadow curriculum problem for
institutional QA: if good practice goes undocumented, internal approval processes
cannot credit or require it.}

\subsection{Bloom's Taxonomy and Verb Observability}

The most widely applied framework for evaluating the cognitive demand of learning
outcomes is Bloom's Taxonomy of Educational Objectives \citep{bloom1956taxonomy} and
its subsequent revision by \citet{anderson2001taxonomy}. The revised taxonomy
distinguishes six levels of cognitive process --- Remember, Understand, Apply, Analyse,
Evaluate, Create --- each associated with a set of observable action verbs. The verb
is the operationalisation of the level: it specifies what the student must \emph{do}
and, by extension, what the assessor must observe.

The critical distinction for quality assurance purposes is between verbs that are
\emph{observable} --- whose instantiation can be directly evidenced in student
performance --- and those that are \emph{unobservable}. Verbs such as
\textit{understand}, \textit{know}, \textit{appreciate}, and \textit{be aware of}
describe internal cognitive states that cannot be directly measured. Their prevalence
in module descriptors signals either weak outcome design or a disconnect between
the language of course documentation and the reality of teaching and assessment
\citep{kennedy2007writing}. Importantly, the problem is not that deep understanding
is an unworthy goal --- it is that the verb provides no specification of what evidence
of understanding would look like.

\subsection{Constructive Alignment}

Biggs's constructive alignment framework \citep{biggs1996enhancing, biggs2011teaching}
provides the theoretical basis for our extension of the analysis beyond verb quality.
Alignment requires internal consistency across three elements: intended learning outcomes,
teaching and learning activities, and assessment tasks. A module in which the LO verbs
are pitched at high cognitive demand (e.g., \textit{design}, \textit{evaluate},
\textit{synthesise}) but the sole assessment instrument is a closed-book formal
examination represents a specific form of misalignment: the assessment cannot elicit
the performance the outcome specifies.

This form of misalignment --- high-demand LOs, low-authenticity assessment ---
is practically significant. It disadvantages students who possess the higher-order
capability described in the outcome but whose performance is measured by a proxy
(recall under examination conditions) that conflates memory with competence. It also
undermines the validity of grades as signals of graduate capability, with consequences
for employers, professional bodies, and the integrity of NFQ level claims
\citep{bloxham2007marking}.

\subsection{Automated Analysis of Learning Outcome Text}

The application of NLP methods to educational text is a growing field
\citep{demetriadis2018, zawacki2019systematic}. Prior work has applied
automated methods to the classification of educational objectives
\citep{khandelwal2021bloom, ullmann2019automated}, the detection of LO quality
issues \citep{chen2022automated}, and cross-programme curriculum mapping
\citep{shabani2012using}. However, corpus-scale cross-institutional analysis
of LO language quality --- spanning an entire national HE system --- has not
previously been reported.

The availability of machine-readable module catalogues across Irish institutions,
combined with advances in rule-based and large language model (LLM) classification,
makes such analysis technically feasible at scale for the first time. Our study
exploits this opportunity.

% ════════════════════════════════════════════════════════════
\section{Research Questions}
\label{sec:rq}
% ════════════════════════════════════════════════════════════

This study addresses three research questions:

\begin{enumerate}[label=\textbf{RQ\arabic*}]
  \item What is the distribution of verb observability tiers across module-level
        learning outcomes in Irish higher education, and does this distribution
        differ systematically between the university and technological university
        sectors?

  \item What is the prevalence of constructive misalignment --- defined as
        high-demand LOs paired with assessment modalities that cannot elicit
        the specified performance --- across the corpus?

  \item What methodological and policy implications follow from corpus-scale
        LO quality analysis for institutional module approval processes and
        national quality assurance frameworks?
\end{enumerate}

% ════════════════════════════════════════════════════════════
\section{Methodology}
\label{sec:methods}
% ════════════════════════════════════════════════════════════

\subsection{Corpus Construction}

Module learning outcome data were collected from sixteen Irish HEIs through systematic
automated retrieval of publicly accessible module catalogue pages during the period
October 2024 -- April 2026. The institutions comprised six public universities
(Maynooth University [MU], University College Dublin [UCD], University College Cork
[UCC], University of Galway [Galway], University of Limerick [UL], and Dublin City
University [DCU]), three technological universities (Atlantic Technological
University [ATU], South East Technological University [SETU], and Munster
Technological University [MTU]), two Institutes of Technology (Dundalk Institute of
Technology [DkIT] and Dún Laoghaire Institute of Art, Design and Technology [IADT]),
three private and specialised colleges (National College of Ireland [NCI], National
College of Art and Design [NCAD], and Griffith College Dublin [GC]), one College of Education (Marino Institute of
Education [MIE]), and TU Dublin [TUD], whose catalogue format is discussed separately
as a methodological finding (Section~\ref{sec:tud}).

\paragraph{Sectoral terminology.}
We use \textit{public universities} (PU) to refer to institutions designated under
the Universities Act 1997. The broader Irish HE system also includes Technological
Universities (TU), Institutes of Technology (IT), Colleges of Education (CE), and
private and specialised colleges (PS). The PU label does not imply that TUs are
private or non-public --- all Irish HEIs receive public funding. The distinction is
legislative and sectoral. \todo{Confirm preferred terminology; check HEA System
Performance Framework for current administrative categories.}

Table~\ref{tab:corpus} summarises the corpus. Each institution required a bespoke
retrieval approach reflecting the diversity of underlying catalogue systems:
CourseLeaf (UCC, Selenium headless Chrome), Akari curriculum management system
(MTU, SETU, DkIT, NCI --- sequential module-ID sweeps), BookOfModules (UL),
the nuigalway.ie course catalogue (Galway), modspec.dcu.ie and dcu.akarisoftware.com
(DCU), TU Dublin Vault (TUD), and institution-specific HTML catalogues for the
remaining institutions. IADT and NCAD module data were extracted from publicly
available PDF catalogues using pdfplumber. MIE data were parsed from four programme
handbooks. In all cases, module-level data were extracted including:
module code, title, ECTS credits, NFQ level where published, and the full text
of each numbered learning outcome.

\begin{table}[H]
\centering
\caption{Corpus composition by institution (D1--D3 scoring complete for all).}
\label{tab:corpus}
\begin{tabular}{llrrrc}
\toprule
\textbf{Institution} & \textbf{Type} & \textbf{Modules} & \textbf{LOs} & \textbf{D1 Mean} & \textbf{Tier 0 (\%)} \\
\midrule
MU      & PU & 2,139  & 12,603 & 1.637 & 10.0 \\
DCU     & PU & 3,464  & 14,472 & 1.682 &  9.5 \\
Galway  & PU & 2,847  & 12,477 & 1.627 &  9.7 \\
UCC     & PU & 6,802  & 29,102 & 1.697 &  4.9 \\
UCD     & PU & 5,419  & 27,145 & 1.158 & 13.4 \\
UL      & PU & 5,731  & 27,056 & 1.736 &  4.9 \\
\midrule
\textit{PU cluster} & & \textit{26,402} & \textit{122,855} & \textit{1.572} & \textit{8.3} \\
\midrule
ATU     & TU & 5,949  & 33,770 & 1.846 & 5.0 \\
SETU    & TU &  3,772 & 18,162 & 1.529 & 1.6 \\
MTU     & TU &  8,217 & 29,881 & 1.846 & 2.8 \\
\midrule
\textit{TU cluster} & & \textit{17,938} & \textit{81,813} & \textit{1.775} & \textit{3.4} \\
\midrule
DkIT    & IT &  4,697 & 19,590 & 1.644 & 6.6 \\
IADT    & IT &    175 &    814 & 1.776 &  4.9 \\
\midrule
NCI     & PS &    749 &  3,647 & 1.627 & 10.4 \\
NCAD    & PS &     79 &    269 & 1.796 &  2.6 \\
GC      & PS &     63 &    441 & 1.914 &  0.5 \\
MIE     & CE &   163 &     49 & 1.592 &  0.0 \\
\midrule
\textbf{Total} & & \textbf{51,190} & \textbf{233,005} & --- & --- \\
\bottomrule
\end{tabular}
\end{table}

\paragraph{ATU: cross-programme deduplication.}
ATU LO data were retrieved from programme pages on \texttt{atu.ie}, where modules
appear in multiple degree programmes. Raw extraction yielded 69,532 LOs across a
stated 14,200 module instances; after removing cross-programme duplicates
(deduplication key: module code $\times$ LO index $\times$ normalised text) the
corpus reduces to 33,770 unique LOs from 5,949 unique module codes --- a 51.4\%
reduction. A further 622 records are retained as shared-code variants (same module
code, different content and credit loading), identified by a version suffix in their
identifier.

\paragraph{DCU: Institute of Education modules.}
The DCU corpus includes modules from the DCU Institute of Education, which
incorporates the former Church of Ireland College of Education (CICE), St Patrick's
College Drumcondra, and Mater Dei Institute (all merged into DCU in 2016).

\paragraph{UL: Mary Immaculate College.}
The UL corpus includes modules validated under UL's degree-awarding authority for
Mary Immaculate College (MIC), Limerick. MIC-specific subject codes (MH, GY, RS) are
confirmed present; B.Ed and PME modules are integrated under shared UL subject codes
and cannot be disaggregated without department-level metadata.

\paragraph{MIE: programme handbook extraction.}
MIE data were extracted from four publicly available programme handbooks (B.Ed.\
Primary, B.Sc.\ ECE, B.Sc.\ Ed.\ Studies, B.Oid.). Of 163 modules across the four
programmes, only 14 (8.6\%) contain explicit student-centred learning outcomes in the
handbook text; the remainder provide module descriptions only. The B.Sc.\ Ed.\ Studies
handbook contains no module-level LO statements. D1 mean and Tier~0 are computed from
the 49 English-language LOs only. The B.Oid.\ (Irish-medium B.Ed.) handbook contains
module descriptions in Irish (Gaeilge); these 38 modules are included in the module
count but excluded from D1 scoring.

\paragraph{Griffith College: PDF extraction.}
GC module LO data were extracted from 74 PDF module descriptor documents retrieved
from Griffith College programme pages. Each PDF follows a standardised format with
module title, ECTS credits, NFQ level, and numbered learning outcomes. 63 of 74 PDFs
contained parseable LO lists; the remainder provided aims-only descriptors and were
excluded. No module codes are published in these PDFs; GC records carry no
\texttt{moduleCode} field.

\paragraph{TU Dublin: dual-source retrieval.}
TU Dublin was surveyed via two routes. The primary public catalogue (Vault system,
5,449 module pages) does not publish student-centred LOs; modules use an aims format
(\textit{``The aim of this module is to [verb]\ldots''}). PDF module catalogues
published by two TUD schools were retrieved separately and yielded 120 modules and
736 student-centred LOs across six programmes (TU835 Architecture, TU874 Mathematical
Sciences, TU078, TU079, TU248, TU5294). The combined TUD corpus (3,527 records) is
included in the totals; the aims-vs-LO distinction is discussed as a methodological
finding in Section~\ref{sec:tud}.

\paragraph{DkIT and IADT.}
Dundalk Institute of Technology (DkIT) contributes 4,697 modules and 19,590 LOs
(D1 mean 1.644). Its D1 mean is notably lower than the other TU-sector institutions
and closer to the university cluster, which may reflect its status as a remaining IoT
rather than a designated TU; this is discussed in Section~\ref{sec:discussion}.
IADT module descriptors were retrieved from a single publicly accessible PDF module
catalogue (175 modules, 814 LOs, D1 mean 1.776). No assessment data are published in
this document; IADT contributes to LO quality analysis only.

\paragraph{Excluded institutions.}
Trinity College Dublin (TCD) is under active investigation but is excluded from the
present analysis due to its decentralised catalogue structure. TUS (Technological
University of the Shannon: Midlands Midwest) and RCSI University of Medicine and
Health Sciences are also not yet included.

\subsection{Module Descriptor Enrichment}
\label{sec:descriptors}

Each LO record in the corpus was augmented with its module descriptor --- the
contextual summary text (typically 50--300 words) that precedes the LO list in
the published module catalogue. Descriptors provide the interpretive context
needed for LLM-based D4--D6 scoring and for the accompanying LO Browser
application (Section~\ref{sec:app}). Descriptor retrieval required institution-specific
strategies, summarised in Table~\ref{tab:descriptors}.

\begin{table}[H]
\centering
\caption{Module descriptor enrichment by institution.}
\label{tab:descriptors}
\begin{tabular}{llrr}
\toprule
\textbf{Institution} & \textbf{Source / Method} & \textbf{Modules} & \textbf{Coverage} \\
\midrule
MU      & Institutional data export (\texttt{moduleContent} field) & 2,537 & 100\% \\
ATU     & \texttt{moduleContent} field in LO records (per-LO text) & 5,941 & $\sim$100\% \\
UCC     & \texttt{ucc-modules-raw.json} description field & 5,137 & $\sim$100\% \\
UL      & \texttt{ul-modules-raw.json} description field & 5,709 & $\sim$100\% \\
MTU     & \texttt{mtu-modules-raw.json} description field & 5,877 & $\sim$96\% \\
TUD     & \texttt{tud-modules-raw.json} description field & 5,385 & $\sim$100\% \\
SETU    & Akari module records description field & 3,628 & $\sim$98\% \\
DkIT    & Akari module records description field & 2,727 & $\sim$92\% \\
UCD     & \texttt{ucd-modules-raw.json} description field & 3,197 & $\sim$62\%$^a$ \\
Galway  & Scraped from \texttt{nuigalway.ie/course-information/} & 1,455 & $\sim$97\% \\
DCU     & modspec.dcu.ie (old and Akari endpoints); dcu.akarisoftware.com & 1,670 & $\sim$48\%$^b$ \\
NCI     & Scraped from \texttt{nci.akarisoftware.com} (indicative content) & 747 & $\sim$100\% \\
NCAD    & Extracted from PDF module descriptors (pdfplumber, Introduction section) & 78 & $\sim$99\% \\
IADT    & Extracted from PDF module catalogue (pdfplumber, aims section) & 112 & $\sim$64\% \\
GC      & Extracted from PDF module descriptors (pdfplumber, Module Objectives section) & 63 & $\sim$100\% \\
MIE     & Not available (programme handbooks contain no descriptor text) & 0 & 0\% \\
\bottomrule
\end{tabular}

\smallskip
\noindent{\footnotesize
$^a$ UCD's catalogue does not publish descriptors for all modules; coverage is limited by
institutional catalogue completeness.\\
$^b$ DCU descriptor coverage reflects an incomplete migration to the Coursebuilder/Akari
system at time of collection; modules not yet migrated have no machine-readable descriptor
in any publicly accessible endpoint.
}
\end{table}

\noindent
Descriptor text was extracted using institution-appropriate methods: direct field
reads from JSON catalogue exports (MU, ATU, UCC, UL, MTU, TUD, SETU, DkIT, UCD);
HTTP scraping with targeted HTML parsing (Galway, NCI, DCU); and PDF text extraction
using \texttt{pdfplumber} with section-boundary detection based on heading patterns
(NCAD: ``1.~Introduction'' section; IADT: aims text preceding the ``On successful
completion'' LO trigger; GC: ``Module Objectives'' or ``Module aims and objectives''
section). Non-English descriptors (primarily Irish-medium MIE and
UL/MIC modules) are retained without language filtering for subsequent analysis.
Descriptor text is stored in the \texttt{moduleContent} field of each normalised LO
record and is available through the LO Browser application.

\subsection{A Tripartite Evaluation Framework: Measurement, Evaluation, and Validation}
\label{sec:tripartite}

A central methodological premise of this study is that LO quality analysis requires
two epistemologically distinct operations that must be kept separate: \textit{measurement}
and \textit{evaluation}. These are not synonyms, and conflating them produces analytically
incoherent results.

\textbf{Measurement} (the NLP pipeline) applies deterministic rules to observable surface
features. It \textit{calculates}: given an LO string, the D1 score is computable without
interpretation, context, or reference to institutional intent. The same algorithm applied
to the same text will always produce the same result. This reproducibility is a strength
for corpus-scale analysis but a limitation for quality assessment: a rule-based system
cannot read an LO against the policy standards it is meant to satisfy, cannot judge
whether a verb choice is appropriate for a given NFQ level and credit load, and cannot
recognise when a syntactically well-formed LO is substantively empty.

\textbf{Evaluation} (the LLM and human raters) applies normative judgement to contextualised
instances. It \textit{assesses}: given an LO, its module context, its institutional
setting, and the policy framework the institution is required to meet, a rater determines
whether the outcome constitutes a quality statement. Evaluation is inherently interpretive.
This is both a limitation (rater effects, anchoring, disciplinary bias) and a capability
(sensitivity to context, policy grounding, and the kind of reasoning that expert quality
reviewers actually apply).

The research question that the LLM and human raters are answering is:

\begin{quote}
\textit{Does this learning outcome, read in context and against the standards its
institution is required to meet, constitute a well-formed, assessable, and
level-appropriate outcome?}
\end{quote}

This is the question a QQI validation panel would ask. It is not reducible to a
calculation.

The third stream --- \textbf{validation via inter-rater reliability (IRR)} --- is the
mechanism by which the measurement and evaluation streams are brought into dialogue.
By computing pairwise agreement between the NLP pipeline, the LLM evaluator, and
expert human raters on a stratified sample, we can determine: (a) the extent to which
automated measurement captures the same variance as expert evaluation; (b) where the
two approaches diverge and why; and (c) whether policy-grounded LLM evaluation achieves
sufficient agreement with human judgement to serve as a scalable quality audit instrument.

\subsection{The D1--D6 Evaluation Rubric}

We developed a six-dimension evaluation rubric (Table~\ref{tab:rubric}) grounded in
the LO quality literature and calibrated to the Irish NFQ context. Each dimension is
scored 0--3, yielding a composite score of 0--18. D1--D3 are operationalised through
the rule-based NLP pipeline (measurement); D4--D6 require contextual judgement and
are operationalised through the LLM evaluator and human raters (evaluation).

\begin{table}[H]
\centering
\caption{The D1--D6 evaluation rubric and the tripartite framework.}
\label{tab:rubric}
\begin{tabularx}{\textwidth}{clXl}
\toprule
\textbf{\#} & \textbf{Dimension} & \textbf{Description} & \textbf{Stream} \\
\midrule
D1 & Verb observability     & Cognitive demand and observability of the leading action verb & NLP (measurement) \\
D2 & Behavioural singularity & Degree to which the LO specifies a single, bounded behaviour & NLP (measurement) \\
D3 & Student-centredness    & Grammatical framing from the student's perspective           & NLP (measurement) \\
D4 & Scope appropriateness  & Whether the scope matches NFQ level and ECTS credit load     & LLM + Human (evaluation) \\
D5 & Assessability signal   & Extent to which the LO implies a specific assessment approach & LLM + Human (evaluation) \\
D6 & NFQ level calibration  & Alignment of cognitive demand to the stated NFQ level        & LLM + Human (evaluation) \\
\bottomrule
\end{tabularx}
\end{table}

For IRR purposes, the LLM evaluator also independently scores D1--D3 on the same
sample, providing a cross-validation signal between the measurement and evaluation
streams on the dimensions where both can operate.

\subsection{NLP Pipeline: Overview and Pre-processing}
\label{sec:nlp-pipeline}

The rule-based NLP pipeline is implemented in a single-pass Node.js scorer
(\texttt{score-rule-based.mjs}) that reads normalised LO records and produces
per-LO scores for D1--D3 together with structured rationale strings and diagnostic
flags. The pipeline is deterministic and reproducible: given the same LO text and
lexicon version, it returns identical output on every run.

\paragraph{Input normalisation.}
Each raw LO text undergoes a cleaning pass before analysis:

\begin{enumerate}
  \item \textit{Numbering prefix removal.} LO text extracted from PDFs and web pages
    frequently retains the source numbering (\texttt{1.}, \texttt{Lo 1.}, \texttt{i.},
    \texttt{(a)}). A family of regular expressions strips these prefixes before
    any linguistic analysis.
  \item \textit{Language detection.} A word-list of high-frequency Irish (Gaeilge)
    function words and verbal particles is checked against the LO text.
    Irish-medium LOs (\texttt{language=ga}) are excluded from D1--D3 scoring and
    assigned a \texttt{null} score with an explanatory flag. This affected 33 modules
    from MIE's B.Oid.\ programme; the B.Oid.\ modules are retained in the corpus
    with their metadata intact.
  \item \textit{Infinitive and framing prefix stripping.}
    The leading verb is isolated by removing common preamble patterns before
    attempting verb extraction: \textit{``To be able to''}, \textit{``To''} (bare
    infinitive), \textit{``Be able to''}, and third-person
    \textit{``Students will (be able to)''} wrappers are stripped so that the true
    action verb is presented for lookup.
\end{enumerate}

\paragraph{Diagnostic flags.}
The pipeline produces a structured \texttt{flags} array per LO for quality audit
purposes. Flags include: \texttt{unknown-verb} (verb not in lexicon, requiring review);
\texttt{no-action-verb} (LO opens with a noun or adjective rather than an action verb,
e.g.\ \textit{``Understanding of\ldots''}, \textit{``Knowledge of\ldots''}); \texttt{wrong-tense}
(past-tense or participial form rather than imperative, e.g.\ \textit{``Demonstrated\ldots''},
\textit{``Applying\ldots''}); \texttt{irish-language}; \texttt{compound} (multiple
behaviours detected); \texttt{instructor-centered}; and \texttt{context-dependent}
(verb whose observability tier is context-sensitive). LOs carrying a null D1 score
are excluded from mean calculations but retained in the corpus.

\subsection{D1 Scoring: Verb Observability Tier}
\label{sec:d1}

D1 operationalises the observability of the leading action verb, grounded in Mager's
(\citeyear{mager1962preparing}) criterion of behavioural observability and the revised
Bloom's taxonomy \citep{anderson2001taxonomy}. The central instrument is a tiered verb
lexicon (v1.3) developed iteratively from the corpus and reviewed against published
LO quality guides from QQI \citeyear{qqi2017validation}, the ECTS Users' Guide
\citeyear{ec2015ects}, and the Cedefop \citeyear{cedefop2017lo} handbook on
defining and writing learning outcomes --- the document explicitly cited by QQI
as the normative reference for LO writing guidance.

\paragraph{Lexicon structure.}
The lexicon (367 entries at v1.3, distributed as supplementary material) assigns
each verb to one of four tiers:

\begin{description}[leftmargin=3em, labelwidth=2.5em]
  \item[Tier 0 — Unobservable] (50 verbs)\quad Internal cognitive states that cannot be directly
    assessed without further specification: \textit{understand, know, appreciate,
    be aware of, learn, grasp, realise, be familiar with, internalise}.
    These map to Mager's \textit{fuzzy} verbs and correspond to Anderson and
    Krathwohl's \textit{Remember} and undifferentiated \textit{Understand} categories
    where no observable performance is implied.
  \item[Tier 1 — Weak observable] (53 verbs)\quad Behaviours that imply an observable
    output but are insufficiently specific to determine the assessment method without
    additional context: \textit{describe, identify, discuss, explain, list, name,
    state, define, outline, summarise}.
  \item[Tier 2 — Clear observable] (207 verbs)\quad Behaviours with a defined
    performance criterion and a clear assessment product: \textit{analyse, evaluate,
    apply, compare, contrast, demonstrate, classify, interpret, solve, justify,
    differentiate, predict}.
  \item[Tier 3 — Precise/contextual] (57 verbs)\quad High-demand behaviours with strong
    domain specificity implying both the action and its expected artefact or product:
    \textit{design, construct, synthesise, critique, formulate, derive, implement
    (novel solution), develop (original), compose, engineer}.
\end{description}

Multi-word verbs (\textit{be aware of}, \textit{critically evaluate}, 18 in total)
are matched using a longest-match greedy strategy before single-token lookup.
Eleven verbs are flagged as \textit{context-dependent}: their observability tier
depends on the domain or the surrounding noun phrase in ways that a lexical lookup
cannot resolve (\textit{understand} is the canonical example, but \textit{reflect},
\textit{investigate}, \textit{communicate}, and \textit{implement} also appear in
this category). Context-dependent verbs are scored by their primary-tier assignment
and flagged for potential LLM re-evaluation in the D6 pass.

\paragraph{Verb extraction algorithm.}
After pre-processing strips leading preamble, the following extraction sequence is applied:

\begin{enumerate}
  \item \textit{Multi-word greedy match.} All multi-word lexicon entries (sorted longest-first)
    are checked; the first match wins.
  \item \textit{``have [past-participle]''} forms. Some LOs use perfect aspect
    (\textit{``have identified\ldots''}); the auxiliary is discarded and the participle
    is used as the verb candidate.
  \item \textit{``be [adjective]''} competency framing. LOs opening with
    \textit{``be [adjective]''} (e.g.\ \textit{``Be competent in\ldots''}) describe
    a state rather than a performance. These receive D1 = 1 (weakly observable state)
    with a \texttt{competency-framing} flag.
  \item \textit{Adverb detection.} The qualifier \textit{critically} is treated as a
    tier-boosting adverb: if the opening token is \textit{critically} and the
    following verb maps to Tier $t < 3$, the score is elevated to $\min(t+1, 3)$.
    Style-only adverbs (\textit{clearly, effectively, confidently}) are stripped
    without affecting the tier.
  \item \textit{Noun/adjective opener detection.} A 32-token set of common noun or
    adjective openers (\textit{understanding, ability, knowledge, skills, competence,
    awareness, \ldots}) triggers a \texttt{no-action-verb} flag and a null D1 score.
    These represent a recurring structural defect in which an LO describes a capability
    attribute rather than an observable performance.
  \item \textit{Wrong-tense detection.} Past-tense and participial forms
    (\textit{developed, demonstrated, applying, exploring}) trigger a
    \texttt{wrong-tense} flag and a null D1 score, with a rationale noting the
    required correction.
  \item \textit{Single-token lookup.} The remaining first token is looked up in the
    verbTier map (case-insensitive). A match returns the tier score; no match returns
    null with an \texttt{unknown-verb} flag.
\end{enumerate}

\paragraph{D1 scoring outcome.}
Each scoreable LO receives D1 $\in \{0, 1, 2, 3\}$. Null scores (Irish-language,
no-action-verb, wrong-tense, unknown-verb) are excluded from institution means
but counted separately. Across the corpus, null D1 scores represent approximately
2.1\% of all LO records.
\todo{Insert final null rate from corpus-wide stats.}

\subsection{D2 Scoring: Behavioural Singularity}
\label{sec:d2}

D2 operationalises the principle that each LO should specify a single bounded,
independently assessable student behaviour \citep{biggs2011teaching, qqi2016core}.
Compound LOs — those that bundle multiple distinct behaviours into a single statement
using coordinating conjunctions — reduce the granularity available for assessment
mapping and progression tracking.

\paragraph{Detection algorithm.}
After stripping the leading infinitive opener (\textit{``To''}), the LO text is
split into clauses at the following boundary tokens: \textit{`` and ''}, \textit{`` or ''},
semicolons, and commas followed by a space and a lowercase letter (to catch comma-listed
verb phrases). For each resulting clause, the first token is extracted and checked
against the verb lexicon. Tokens not in the lexicon are not counted. The count of
distinct recognised verbs (\textit{not} the raw clause count) determines the D2 score:

\begin{center}
\begin{tabular}{cl}
\toprule
\textbf{D2 score} & \textbf{Criterion} \\
\midrule
3 & Single behaviour — one recognised action verb, no coordination \\
2 & Two distinct behaviours — compound but manageable \\
1 & Three behaviours — conflated; splitting recommended \\
0 & Four or more behaviours — highly fragmented \\
\bottomrule
\end{tabular}
\end{center}

\paragraph{Limitations.}
The detection is conservative by design: only tokens present in the verb lexicon are
counted, so a compound LO built from unlexiconed verbs may be under-counted.
Adversely, a single verb with several qualifying noun phrases (e.g., \textit{``analyse
the structure, function, and regulation of\ldots''}) is not falsely scored as compound.
D2 is best interpreted as a lower bound on behavioural multiplicity; the LLM pass
(D5, assessability signal) is intended to provide a context-sensitive complement.

\subsection{D3 Scoring: Student-Centredness}
\label{sec:d3}

D3 assesses the grammatical framing of the LO, distinguishing LOs that place the
student's observable performance at the centre from those framed from the module's
or instructor's perspective. The distinction is required by QQI policy
\citeyear{qqi2016core} and the Bologna/ECTS framework \citeyear{ec2015ects}, both
of which specify that LOs must describe what the \textit{student} will be able
to demonstrate upon completion, not what the \textit{module} will cover.

\paragraph{Scoring scheme.}
Patterns are checked in priority order (first match wins):

\begin{description}[leftmargin=3em, labelwidth=2.5em]
  \item[D3 = 0 — Instructor/module-centred] LO begins with one of 14 instructor-centred
    patterns: \textit{``This module will\ldots''}, \textit{``The aim of this module
    is to\ldots''}, \textit{``This course will introduce\ldots''}, etc.
    These are aims statements masquerading as LOs; the grammatical subject is the
    module or course, not the student.
  \item[D3 = 1 — Passive student] LO places the student as a recipient rather than an
    agent: \textit{``Students will be introduced to\ldots''}, \textit{``Students will
    be provided with\ldots''}, \textit{``Students will become familiar with\ldots''}.
    Nine such passive constructions are detected. These are student-centred in framing
    but describe an input to the student rather than an output from the student.
  \item[D3 = 2 — Active student, indirect framing] LO states what the student will do,
    but uses a third-person or infinitive wrapper: \textit{``Students will [verb]\ldots''},
    \textit{``The student will be able to\ldots''}, \textit{``To [verb]\ldots''}.
    These are student-centred and assessable but use a more indirect grammatical
    structure than the preferred OBE convention.
  \item[D3 = 3 — Verb-first imperative] LO opens directly with a recognised action verb,
    placing the student performance as the grammatical subject by implication:
    \textit{``Analyse\ldots''}, \textit{``Design\ldots''}, \textit{``Demonstrate\ldots''}.
    This is the preferred convention in outcome-based education
    \citep{biggs2011teaching} and is explicitly recommended by QQI \citeyear{qqi2016core}.
\end{description}

\paragraph{Relationship to D1.}
D3 and D1 are orthogonal by construction: a verb-first LO can score D3 = 3 with a
Tier 0 verb (\textit{``Understand the principles of\ldots''}), and a module-centred
framing can contain a high-tier verb (\textit{``This module will develop students'
ability to synthesise\ldots''}). The two dimensions capture independent quality
properties and are reported separately throughout the analysis.

\subsection{LLM Evaluation: D4--D6 (Pilot)}
\label{sec:llm-eval}

The three rule-based dimensions (D1--D3) operate on lexical and syntactic surface
features and are context-free by design: the D1 score for \textit{``understand the
principles of machine learning''} is the same regardless of whether the module is
a first-year introductory unit or a Level 9 specialist course. The remaining three
dimensions (D4 scope appropriateness, D5 assessability signal, D6 NFQ calibration)
require contextual judgement that is beyond the capacity of rule-based methods,
and are therefore operationalised through an LLM evaluator.

\paragraph{Architecture.}
The evaluator (\texttt{score-llm-claude.mjs}) makes a single structured API call
per LO to the Anthropic Claude API, using a prompt that combines:

\begin{enumerate}
  \item \textbf{Module context:} module code, title, department, ECTS credits, NFQ level,
    and the full module descriptor text (truncated to 200 words where necessary).
  \item \textbf{The LO under evaluation,} identified by its within-module index.
  \item \textbf{The D4--D6 rubric,} presented as a structured scoring instrument with
    explicit 0--3 anchor descriptions per dimension.
  \item \textbf{NFQ level descriptors,} mapping each integer level to the QQI knowledge,
    skill, and competence expectations for that level.
\end{enumerate}

The model is instructed to respond with valid JSON in a fixed schema
(\texttt{\{d4: \{score, rationale\}, d5: \{score, rationale\}, d6: \{score, rationale\}\}})
and to reason step by step before committing to each score. Response temperature is
left at the API default (not set); scores are clamped to $[0, 3]$ and rounded to the
nearest integer. Responses that do not parse as valid JSON, or that are missing any
required dimension, are logged as errors and excluded from analysis.

\paragraph{Current status.}
A pilot evaluation has been completed on the MU corpus
(\todo{insert $n$ from scored-llm-claude.json}), using \texttt{claude-haiku-4-5}
as the evaluating model. Full cross-institutional evaluation is planned but not yet
complete. Results from the pilot are not reported in the current paper; this section
documents the methodology for the follow-on analysis.

\paragraph{Policy-grounded prompt architecture.}
The v2 evaluator (\texttt{score-llm-v2.mjs}) extends the pilot with a two-tier prompt
structure designed to ground the LLM in the same normative framework as an expert
human rater. The \textbf{system prompt} (fixed across all calls, eligible for Anthropic
prompt caching) contains:

\begin{itemize}
  \item \textbf{QQI Policies and Criteria for the Validation of Programmes}
    \citep{qqi2017validation} and \textbf{QQI Core Policies and Criteria}
    \citep{qqi2024core} — the Irish national standards requiring that module LOs
    are explicitly specified, plainly written, not vague or generalised, and
    appropriate to NFQ level and credit volume.
  \item \textbf{QQI Assessment and Standards} \citep{qqi2022assessment} — requiring
    assessment to be criterion-referenced against stated outcomes.
  \item \textbf{ECTS Users' Guide} \citep{ec2015ects} — the Bologna framework
    requiring demonstrable achievement and discouraging unobservable verbs such as
    \textit{understand} and \textit{know}.
  \item \textbf{Cedefop (2017)} \citep{cedefop2017lo} — QQI's stated normative
    reference for LO writing guidance, including active verb requirements, specificity
    criteria, and measurability standards.
  \item \textbf{NFQ level descriptors} (levels 6--10) — knowledge, know-how/skill,
    and competence expectations at each level, used to ground D6 calibration.
  \item \textbf{The D1--D6 rubric}, presented as a structured scoring instrument
    with explicit 0--3 anchor descriptions per dimension.
\end{itemize}

The \textbf{user prompt} (per LO) contains an institution-specific context block
selected by institution code, drawing from a corpus of LO writing guidance documents
collected from all sixteen institutions in the corpus. These documents range from
full standalone institutional guides (DCU Teaching Enhancement Unit, 2023; UCC CIRTL,
2023) to QA framework appendices and curriculum review checklists. Where an institution
has adopted QQI or Kennedy \citeyear{kennedy2007writing} as its primary reference without
supplementary guidance, the evaluator falls back to the shared policy layer. The
user prompt additionally provides module code, title, ECTS credits, NFQ level, and
module descriptor text.

This architecture means the evaluator is answering the question a QQI validation
panel would ask: not merely whether an LO contains an action verb, but whether it
meets the standards the institution is itself required to apply.

\paragraph{Computational efficiency: prompt amortisation through batching.}
A policy-grounded LLM evaluator introduces a structural cost asymmetry: the system
prompt --- containing the full QQI, ECTS, Cedefop, and NFQ policy layer plus the
evaluation rubric --- is approximately 1,200 tokens and is fixed per API call, while
the per-LO user content (institutional context block, module descriptor, and LO text)
is approximately 265 tokens. When LOs are evaluated one per call, the fixed cost
dominates: approximately 82\% of input tokens are system prompt, applied to a single
LO. Sending $N$ LOs per call amortises the system prompt across $N$ instances,
reducing the effective input token cost per LO to:

\begin{equation}
  C_{\text{input}}(N) = \frac{T_{\text{sys}}}{N} + T_{\text{lo}}
\end{equation}

where $T_{\text{sys}}$ is the system prompt token count and $T_{\text{lo}}$ is the
per-LO user content token count.

Table~\ref{tab:batch} shows empirical results from the MU pilot
(\texttt{gpt-4o-mini}; $n = 50$ per condition; 300 output tokens budgeted per LO).

\begin{table}[H]
\centering
\caption{Effect of batch size on LLM evaluation cost and throughput (MU pilot, gpt-4o-mini).}
\label{tab:batch}
\begin{tabular}{rrrrrr}
\toprule
\textbf{Batch} & \textbf{Input tok/LO} & \textbf{Output tok/LO} & \textbf{Cost/LO} & \textbf{LOs/min} & \textbf{Full corpus est.} \\
\midrule
1  & 1,465 & 162 & \$0.000317 & 14.6 & \$73.9 \\
5  &   516 & 163 & \$0.000175 & 16.6 & \$40.8 \\
10 &   398 & 143 & \$0.000146 & 31.3 & \$34.0 \\
\bottomrule
\end{tabular}
\end{table}

The input token saving follows the expected curve: moving from batch~1 to batch~5
reduces per-LO input tokens by 65\%; moving from batch~5 to batch~10 yields a further
23\% reduction, reflecting the diminishing-returns character of the harmonic series.
Output tokens are essentially constant across batch sizes (143--163 tokens per LO),
since each LO's evaluation is generated independently within the response. Throughput
at batch~10 doubles relative to batch~1 (31.3 vs 14.6 LOs/min) because fewer API
round-trips are required regardless of per-call latency.

Error rates were zero at both batch~5 and batch~10 when output token budget was set to
$300 \times N$ tokens per call. A tighter budget (150 tokens per LO, used in initial
testing) produced truncation errors on LOs with long module descriptors; increasing
the budget resolved this without materially increasing cost, since output tokens for
short responses are not charged for unused budget.

The practical recommendation for corpus-scale evaluation is batch~5 as a conservative
choice (well-tested, low error rate, 45\% cost reduction vs single-LO) or batch~10
for maximum efficiency (54\% cost reduction, doubled throughput, zero errors in
pilot). Batches larger than~10 are not recommended: the model must track more
simultaneous evaluation targets within a single response, increasing the risk of
cross-LO confusion, and the marginal cost saving beyond batch~10 is small.

This batching approach generalises to any policy-grounded LLM evaluation task where
the system prompt is substantially larger than the per-instance user content --- a
common configuration when providing regulatory or institutional context for document
quality evaluation.

\paragraph{Inter-rater reliability design.}
A stratified sample of 300--500 LOs will be drawn from across the corpus (stratified
by institution sector, NFQ level, and D1 tier to ensure coverage of edge cases) and
evaluated by three raters: (1) the rule-based NLP pipeline (D1--D3); (2) the LLM
evaluator (D1--D6, providing an independent D1--D3 signal for cross-validation);
and (3) one or more expert human raters. Pairwise inter-rater reliability will be
reported as Krippendorff's $\alpha$ for ordinal D1--D6 scales and Cohen's $\kappa$
for D3 as a binary (student-centred / not student-centred) measure.
The IRR analysis is reported in a companion paper.
\todo{Update when IRR data are available.}

\subsection{Human Evaluation}
\label{sec:human-eval}

Human evaluation provides the interpretive anchor against which both the NLP
and LLM streams are validated. Expert raters --- academic staff with experience in
programme design, QQI validation processes, or quality assurance --- apply the
D1--D6 rubric independently to a stratified sample of the corpus, with the same
policy context documents available to the LLM evaluator.

Human evaluation is not positioned as a superior method but as a different one.
It brings disciplinary sensitivity and institutional knowledge that neither the NLP
pipeline nor the LLM evaluator can fully replicate; it is also subject to rater
fatigue, anchoring effects, and disciplinary bias that automated methods avoid.
The three streams are therefore complementary: the NLP pipeline provides scalable,
reproducible measurement across the full corpus; the LLM evaluator provides
policy-grounded evaluation at scale; and the human raters provide the interpretive
anchor that validates both.

\paragraph{Rating tool.}
To facilitate distributed human evaluation, a browser-based rating interface was
developed and deployed at
\url{https://jkeatingmu.github.io/MU-finder/slo-study/}
(see also Appendix~\ref{app:rating-tool}). The tool is fully static --- it requires
no login, no plugin, and no software installation --- and is accessible to any
collaborator with the URL. Raters browse the full 233,005-LO corpus by institution,
expand any LO record to view the full module context and NLP/AI evaluation scores,
and submit ratings via an in-app modal.

The modal presents the LO text and module context at the top for reference, then
solicits scores (0--3) and optional rationale text for each of the four human-stream
dimensions: D1 (verb observability, for IRR against the NLP measurement), D4 (scope),
D5 (assessability), and D6 (NFQ calibration). Scores are submitted anonymously to a
cloud database (Supabase); no personal identifying data are collected or transmitted.
The rater's browser records which LOs have been rated in local storage, allowing the
interface to indicate previously rated records and warn of potential duplicate
submissions; this record is local only and is never transmitted.

\paragraph{IRR sample design.}
Inter-rater reliability will be computed on a stratified sample of approximately
200--300 LOs, drawn proportionally by institution type (university, technological
university, institute), NFQ level (6--10), and D1 tier (0--3) to ensure coverage
of the quality gradient. Each LO in the sample will be independently rated by at
least two human raters. Agreement will be assessed using Cohen's weighted $\kappa$
for each dimension, with a target $\kappa \geq 0.70$ as the threshold for acceptable
reliability \citep{landis1977measurement}. The same sample will be evaluated by the
LLM pipeline, enabling a three-way comparison: NLP measurement vs.\ LLM evaluation
vs.\ human evaluation.

\paragraph{Rater protocol.}
Raters will receive a briefing document presenting the rubric, the underlying policy
framework (QQI QP.17, ECTS Users' Guide, Cedefop 2017, NFQ descriptors), and ten
calibration LOs with consensus scores and written rationale. Raters are instructed
to evaluate each LO as an independent unit --- without anchoring to other outcomes
within the same module --- and to apply the rubric as written rather than adjusting
for perceived institutional norms.

\todo{Human evaluation data collection in progress --- IRR results to be reported
when minimum sample threshold is reached.}

\subsection{Constructive Alignment Analysis}

For the constructive alignment extension, we matched module D1 scores against
assessment modality data. At Maynooth University, assessment data were extracted
from the institutional module descriptor system: each module records the percentage
weight of continuous assessment (CA) versus formal examination. At MTU, the Akari
assessment management system provides richer data, including per-component
assessment type, weight, and --- critically --- a mapping from each component to the
specific LO indices it is designed to assess.

Misalignment was operationalised as follows: a module is classified as \textbf{Type A}
(confirmed misalignment) when its mean D1 verb tier $\geq 2.0$ (clear observable or
above), its formal examination weight $\geq 70\%$, and (for MTU) the high-demand LOs
are specifically mapped to the examination component with no CA rescue. A module is
classified as \textbf{partially rescued} when high-demand LOs appear in both the
examination and CA component mappings. Type B misalignment (low D1 + high CA weight)
captures the converse problem: assessment practice that exceeds the aspiration of
the LO language.

% ════════════════════════════════════════════════════════════
\section{Results}
\label{sec:results}
% ════════════════════════════════════════════════════════════

\subsection{Cross-Institutional D1 Distribution}

Figure~\ref{fig:d1dist} \todo{Insert bar chart: D1 mean by institution, coloured by
sector (PU/TU/PS)} shows the D1 mean for each of the ten institutions. The
within-PU-sector consistency is notable, with five of six universities returning means
between 1.627 and 1.736; UCD (1.158) is an outlier requiring methodological scrutiny
(see below). TU means range from 1.529 (SETU) to 1.846 (ATU and MTU), with SETU
occupying a position closer to the university cluster than the other TUs --- a finding
discussed in Section~\ref{sec:discussion}. The sector gap of $\Delta = 0.203$ (PU mean
1.572, TU mean 1.775) is consistent in direction across every institution; no TU mean
falls below the five-university mean excluding UCD (1.695).

\begin{table}[H]
\centering
\caption{D1 tier distribution by sector.}
\label{tab:tiers}
\begin{tabular}{lrrrrrr}
\toprule
& \multicolumn{2}{c}{\textbf{Tier 0}} & \multicolumn{2}{c}{\textbf{Tier 1}} & \multicolumn{2}{c}{\textbf{Tier 2+}} \\
\cmidrule(lr){2-3} \cmidrule(lr){4-5} \cmidrule(lr){6-7}
\textbf{Sector} & \textbf{n} & \textbf{\%} & \textbf{n} & \textbf{\%} & \textbf{n} & \textbf{\%} \\
\midrule
University  & 10,198 &  8.3 & 71,256 & 58.0 & 41,401 & 33.7 \\
TU          &  3,827 &  3.7 & 51,243 & 49.5 & 48,370 & 46.8 \\
\bottomrule
\end{tabular}
\end{table}

The Tier 0 rate differential is particularly notable: 8.3\% of PU LOs
open with an unobservable verb versus 3.4\% in the TU sector ($\Delta = 4.9$
percentage points). At UCD, which returns the lowest D1 mean in the corpus (1.158),
13.4\% of LOs open with a Tier 0 verb. The UCD D1 mean is substantially below
the five-university mean excluding UCD (1.695), which may reflect formatting
artefacts in UCD's module descriptor export --- specifically, a 31.2\% rate of
null or near-null LO fields --- as well as genuine quality variation; this is
flagged for follow-up in the IRR phase (D4--D6 LLM scoring). The dominant
Tier 0 verb across all universities is \textit{understand} (appearing in 2,065
UCD LOs alone); at TUs, \textit{understand} is present but is not the dominant
Tier 0 form.

The TU advantage on verb observability may reflect the vocational and professional
orientation of TU programmes, in which the performative specification of outcomes
--- what the graduate can \emph{do in a professional context} --- is more deeply
embedded in curriculum culture \citep{dait2022tu}. It may also reflect the
influence of QQI programme validation processes on TU module design, or historical
differences in how Bloom's taxonomy has been adopted and enforced in curriculum
approval workflows across sector types.

\subsection{Dominant Tier 0 Verbs}

Table~\ref{tab:tier0verbs} lists the fifteen most frequent Tier 0 verbs across
the full corpus. The dominance of \textit{understand} is consistent with prior
findings in the international literature \citep{kennedy2007writing}; it is
sufficiently common to suggest institutional habit rather than individual authorial
choice.

\begin{table}[H]
\centering
\caption{Top 15 Tier 0 verbs (full corpus, $n$ = 226,295 LOs).}
\label{tab:tier0verbs}
\begin{tabular}{lrr}
\toprule
\textbf{Verb} & \textbf{Count} & \textbf{\% of Tier 0} \\
\midrule
understand        & \todo{n} & \todo{\%} \\
appreciate        & \todo{n} & \todo{\%} \\
know              & \todo{n} & \todo{\%} \\
be aware of       & \todo{n} & \todo{\%} \\
learn             & \todo{n} & \todo{\%} \\
gain              & \todo{n} & \todo{\%} \\
be familiar with  & \todo{n} & \todo{\%} \\
comprehend        & \todo{n} & \todo{\%} \\
recognise         & \todo{n} & \todo{\%} \\
grasp             & \todo{n} & \todo{\%} \\
explore           & \todo{n} & \todo{\%} \\
familiarise       & \todo{n} & \todo{\%} \\
enhance           & \todo{n} & \todo{\%} \\
cover             & \todo{n} & \todo{\%} \\
develop awareness & \todo{n} & \todo{\%} \\
\bottomrule
\end{tabular}
\end{table}

\note{Run verb frequency query across full corpus to populate Table 3.}

\subsection{D2 and D3: Behavioural Singularity and Student-Centredness}

Across the full corpus, the majority of LOs demonstrate adequate student-centredness
(D3 $\geq 2$; \todo{n}\%), reflecting widespread adoption of the standard
\textit{``On completion of this module the student will be able to\ldots''} framing.
However, behavioural singularity (D2) reveals meaningful variation: \todo{n}\% of LOs
contain two or more distinct action verbs in a single statement, and \todo{n}\% contain
three or more. Compound LOs of this kind resist reliable assessment design and create
ambiguity about which competence is being evaluated.

\subsection{Institutional Format: The TU Dublin Finding}
\label{sec:tud}

TU Dublin presented a methodologically distinct case. Scraping of 5,449 publicly
accessible module catalogue pages (TUD's \textit{Vault} system) revealed that TUD
does not publish student-centred learning outcomes on its public module pages. Instead,
modules are described using a module aims format: \textit{``The aim of this module
is to [verb]\ldots''}. This format is grammatically and epistemologically distinct
from an LO: it describes the module's instructional intention rather than the
student's expected achievement.

Of 5,449 pages queried, 2,687 (49.3\%) yielded extractable aim statements; 89
contained explicit \textit{``students will be able to\ldots''} LOs. The mean D1
score for extracted TUD aims (1.458) is lower than that of any other TU in the
corpus and is comparable to the university cluster. However, direct comparison
with institutions publishing student-centred LOs is methodologically inappropriate;
the aims format is itself the finding.

A supplementary investigation discovered that TU Dublin does publish full
QQI-compliant module catalogues in PDF form for at least one school. Four PDF
module catalogues for the \textit{Spatial Planning and Environmental Management}
programme (TU835, Stages 1--4, 2024--25) were retrieved from the TU Dublin public
website and parsed using a \texttt{pdfplumber}-based extractor targeting the Akari
M1 template format. Subsequent discovery identified a second PDF format --- TU Dublin's
Coursebuilder (Akari M2) template --- used for five further programmes. Systematic
URL probing of all 26 TUD school media directories against 542 programme codes
(147 from the CAO handbook plus codes identified via the TU Dublin A-Z courses page
and a supplementary sweep of the TU001--TU199 short-code range)
yielded PDFs for six programmes in total:

\begin{itemize}
  \item \textbf{TU835} (Architecture/Spatial Planning, 4 stages): 35 modules, 206 LOs in M1 template format
  \item \textbf{TU874} (Mathematical Sciences, 4 stages): 61 modules, 362 LOs in Coursebuilder format
  \item \textbf{TU078} (Mathematical Sciences postgraduate, 1 stage): 8 modules, 59 LOs in Coursebuilder format
  \item \textbf{TU079} (Mathematical Sciences postgraduate, 1 stage): 8 modules, 56 LOs in Coursebuilder format
  \item \textbf{TU248}: 4 modules, 27 LOs in Coursebuilder format
  \item \textbf{TU5294}: 4 modules, 26 LOs in Coursebuilder format
\end{itemize}

Combined, the PDF corpus yields 120 unique modules and 736 student-centred LOs
with a mean D1 score of 1.837 --- substantially higher than the Vault aims corpus
and consistent with TU-sector norms. These records are included in the corpus.
The remaining TUD schools do not appear to have published equivalent PDFs through
discoverable URL patterns; contact with the TUD Quality Office for comprehensive
coverage is recommended for future work. TUD is, to our knowledge,
the only Irish HEI whose primary public module catalogue does not publish
student-centred learning outcomes system-wide, representing a transparency concern
independent of what may exist in internal programme documentation.

\subsection{Specialist Sector: NCAD}

The National College of Art and Design (NCAD) presents a methodologically and
substantively distinctive case. Module descriptors were retrieved from NCAD's public
website as individual PDF files and parsed to extract LO text, assessment components,
and LO-to-assessment mappings. The resulting corpus (79 modules, 269 LOs) is small
relative to institutional HEI corpora but constitutes the complete publicly available
set for the undergraduate programmes surveyed.

NCAD returns the highest D1 mean in the corpus (1.796), with Tier 0 at 2.6\% ---
lower than any other institution. This is consistent with the performative nature of
art and design education: module LOs are characterised by observable, practice-based
verbs (\textit{demonstrate, design, develop, create, produce, exhibit}) that map
naturally to high D1 tiers. The absence of Tier 0 language (\textit{understand,
appreciate}) is itself a disciplinary marker.

Equally notable is NCAD's assessment profile: across all 79 modules, zero formal
examinations are recorded. Assessment is exclusively portfolio-, project-, and
presentation-based --- a complete CA model. This represents a natural point of
comparison with the confirmed Type A modules identified at MTU, SETU, and NCI,
where high-demand LOs are paired with examination-heavy assessment; at NCAD,
the assessment modality is structurally aligned to the LO language by institutional
design.

\subsection{Private Sector: Griffith College Dublin}

Griffith College Dublin (GC) is a designated higher education institution operating
under the Qualifications and Quality Assurance (Education and Training) Act 2012.
Module LO data were retrieved from PDF module descriptor documents linked from
programme pages on the Griffith College website. Seventy-four unique PDFs were
identified and downloaded; 63 (85.1\%) contained parseable numbered learning
outcome lists. The remaining 11 PDFs provided aims-only or empty descriptor text
and were excluded.

PDFs follow a standardised format that includes module title, ECTS credits, NFQ level,
numbered LOs (using Arabic or Roman numerals across two format variants), and a
``Module Objectives'' or ``Module aims and objectives'' section used as the
descriptor field. No module codes are published; GC records carry the module title
as the primary identifier.

The 63 modules yield 441 student-centred LOs --- a mean of 7.0 LOs per module,
higher than the corpus-wide average. GC returns the highest D1 mean in the corpus
(1.914), with Tier 0 at 0.5\% (2 LOs), the second-lowest Tier 0 rate after MIE.
The verb profile is markedly action-oriented; modal verbs (\textit{apply, analyse,
evaluate, design, implement}) dominate, consistent with a strong professional and
applied orientation across Griffith's programme portfolio. The distribution of
assessment data was not available in the PDF format and GC contributes to LO quality
analysis only.

\subsection{College of Education Sector: MIE}

Marino Institute of Education (MIE) contributes the smallest and methodologically
most distinctive sub-corpus in this study. Four undergraduate programme handbooks
were parsed: B.Ed.\ (Primary), B.Sc.\ Early Childhood Education, B.Sc.\ Education
Studies, and B.Oid.\ (Irish-medium B.Ed.). The B.Oid.\ handbook is in Irish (Gaeilge)
and is excluded from D1 scoring, as the verb lexicon is English-only; its modules
are retained in the corpus with a \texttt{language=ga} flag.

Of 163 modules across the four programmes, only 14 (8.6\%) contain student-centred
learning outcomes meeting the trigger criterion (\textit{On/Upon successful completion
of this/the module}, or the Irish equivalent \textit{Ar chríochnú an mhodúil});
the remaining 149 modules (91.4\%) carry module descriptions but no formal LOs.
B.Sc.\ Education Studies presents zero LO-bearing modules across its entire
programme. The 49 scoreable LOs yield a D1 mean of 1.592 with Tier 0 at 0.0\%.

The headline LO-absence rate is less informative than the within-institution
pattern: modules \textit{within the same programme and academic year} differ in
whether they carry formal LOs, with no apparent institutional standard governing
which module types are required to do so. This is not a case of a sector-wide
convention against LO publication (such as might partially explain TUD's
aims-format corpus); it is an inconsistency of practice within a single institution.
For a teacher education college whose graduates will be expected to write LOs for
their own classroom teaching, the absence of a consistent institutional LO
framework is itself a notable finding. Tier 0 at 0.0\% --- the lowest in the
corpus --- suggests that where LOs \textit{are} written at MIE, verb choice is
consistently observational; the quality deficit lies in coverage, not in
the language used where LOs exist.

\subsection{Constructive Alignment: Maynooth University}

At MU, 2,139 modules had complete D1 scoring and published assessment weight data.
Cross-tabulating D1 mean tier against examination weight revealed 73 modules with
mean D1 $\geq 2.0$ and examination weight $\geq 70\%$. Applying a third filter ---
zero reported laboratory or practical hours in the timetabled T\&L activities ---
reduced this to \textbf{49 confirmed Type A} modules: modules in which high-demand
LOs are promised, examination-heavy assessment is applied, and no practical learning
activity is recorded that might provide an alternative alignment pathway.

The remaining 24 modules (high D1, exam $\geq 70\%$, but with laboratory hours
present) were classified as ambiguous: the CA component might rescue alignment
if the CA is authentically aligned to the high-demand LOs, but MU's data recording
practice --- CA as a single aggregate percentage weight, without component type ---
does not permit this determination from data alone.

\subsection{Constructive Alignment: Munster Technological University}

MTU's Akari assessment management system enabled a substantially more precise analysis.
Of 4,987 modules with both D1 scoring and confirmed assessment data, \textbf{114}
met the Type A threshold (mean D1 $\geq 2.0$, exam $\geq 70\%$). The LO-to-component
index in Akari maps each assessment component to the specific LO indices it is
designed to assess, enabling the following classification:

\begin{itemize}
  \item \textbf{55 confirmed misaligned}: high-demand LOs (D1 $\geq 2$) are
    specifically mapped to the formal examination component, with no CA component
    also targeting those LOs. The examination is the sole instrument for assessing
    the highest-demand outcomes.

  \item \textbf{101 partially rescued}: the module is examination-heavy overall, but
    high-demand LOs are also mapped to at least one CA component. The headline exam
    percentage is misleading as a misalignment indicator; the internal LO mapping
    reveals adequate alignment.

  \item \textbf{1 fully aligned}: all high-demand LOs are mapped exclusively to CA
    components despite the high examination weight.

  \item \textbf{7}: LO index data insufficient to classify.
\end{itemize}

The distribution of CA types in confirmed misaligned MTU modules is revealing
(Table~\ref{tab:catypes}). Short Answer Questions (SAQ) is the most frequent CA
type (20 instances), which is effectively examination-like in cognitive structure:
it assesses recall and short-form reproduction rather than the higher-order
performance implied by the high-demand LO verbs.

\begin{table}[H]
\centering
\caption{CA component types in MTU confirmed misaligned modules ($n$ = 55).}
\label{tab:catypes}
\begin{tabular}{lr}
\toprule
\textbf{CA Component Type} & \textbf{Count} \\
\midrule
Short Answer Questions     & 20 \\
Other                      & 14 \\
Written Report             & 12 \\
Project                    & 12 \\
Practical/Skills Evaluation & 10 \\
Essay                      &  7 \\
Open-book Examination      &  5 \\
Presentation               &  4 \\
Non-MTU Exam               &  1 \\
\bottomrule
\end{tabular}
\end{table}

\subsection{Constructive Alignment: South East Technological University}

SETU's Akari data provides a third instance of the gold-standard LO-to-component
analysis. Of 3,772 modules with D1 scoring and confirmed assessment data, \textbf{7}
meet the confirmed Type A threshold (mean D1 $\geq 2.0$, exam $\geq 70\%$, high-demand
LOs mapped exclusively to the examination component), with a further \textbf{40}
partially rescued. SETU also yields \textbf{5 confirmed Type B} modules --- low D1
(mean $< 1.5$) paired with CA-heavy assessment --- representing the converse alignment
problem: assessment that exceeds the aspiration of the LO language. The CA-only
component type in SETU confirmed misaligned modules is uniformly Assignment,
consistent with the observation that written assignments are the most common CA
rescue mechanism across Akari institutions.

The lower absolute count of Type A modules at SETU relative to MTU (7 vs 55) is
partially attributable to SETU's lower overall D1 mean (1.529 vs 1.846): fewer modules
clear the D1 $\geq 2.0$ threshold. This raises the question of whether a lower
Type A count represents better alignment or simply lower LO ambition.

\subsection{Constructive Alignment: National College of Ireland}

NCI's publicly accessible Akari instance provides a fourth gold-standard dataset.
Of 749 modules analysed, \textbf{2} are confirmed Type A and \textbf{26} partially
rescued. No Type B modules were identified. NCI's D1 mean (1.627) is comparable to
the university cluster, despite NCI's classification as a private/specialised
institution; its Tier 0 rate (10.4\%) is the highest in the non-university sample.

Table~\ref{tab:comparison} summarises the alignment findings across all four institutions.

\begin{table}[H]
\centering
\caption{Constructive alignment findings by institution and data method.}
\label{tab:comparison}
\begin{tabular}{lllll}
\toprule
& \textbf{MU} & \textbf{MTU} & \textbf{SETU} & \textbf{NCI} \\
\midrule
Modules with D1 + assessment data & 2,139 & 4,987 & 3,772 & 749 \\
Type A candidates (D1 $\geq 2$, exam $\geq 70\%$) & 73 & 114 & 47 & 28 \\
Confirmed misaligned & 49$^a$ & 55$^b$ & 7$^b$ & 2$^b$ \\
Partially rescued & 24$^c$ & 101 & 40 & 26 \\
Type B (low D1, CA-heavy) & n/a & --- & 5 & 0 \\
Method & Proxy & LO-index & LO-index & LO-index \\
CA type data available & No & Yes & Yes & Yes \\
\bottomrule
\end{tabular}
\end{table}

\noindent
$^a$ Confirmed by lab-hours proxy (zero practical hours). $^b$ Confirmed by
LO-to-component index mapping. $^c$ Ambiguous: exam $\geq 70\%$ but lab hours present; CA may rescue alignment.

The methodological contrast between MU (weight-only + lab-hours proxy) and the three
Akari institutions (MTU, SETU, NCI) is itself a substantive finding.
The 24 MU modules classified as ambiguous would be resolvable with component-level
CA recording. The 101 MTU modules that appear misaligned on exam-weight alone but
are resolved as partially aligned by the LO index are modules that a weight-only
analysis would misclassify. The precision of the Akari data is a direct consequence
of the system's design; it suggests that the data infrastructure of assessment
management systems has a direct bearing on the quality of evidence available for
institutional review. The confirmed misalignment count across all four institutions
(113 confirmed Type A modules) represents a conservative lower bound; the five
institutions with weight-only or no assessment data would yield additional cases
under the same criteria.

\subsection{Readability of Module Descriptors and Learning Outcomes}
\label{sec:readability}

\note{This section reports planned readability analysis on module descriptor text and
LO statements across the full corpus. No API cost — pure text processing.}

A learning outcome that specifies an observable behaviour in the correct grammatical
form may nonetheless be inaccessible to the student it is written for. Verb
observability (D1) captures the \emph{epistemic} quality of an LO statement; its
\emph{communicative} accessibility is a distinct dimension. Module descriptors and
LO text are public-facing documents: students consult them when choosing programmes,
planning their studies, and understanding what will be assessed. Their readability
is therefore a student equity concern as much as a quality assurance one.

We apply a suite of established readability measures to module descriptor text and
LO statements across all institutions in the corpus:

\begin{itemize}
  \item \textbf{Flesch-Kincaid Grade Level (FKGL)}: estimates the US school grade
    level required to comprehend a text; based on sentence length and syllable count.
    A FKGL of 12--14 is typical for HE module descriptors; values above 16 suggest
    prose that may impede student comprehension.

  \item \textbf{Flesch Reading Ease (FRE)}: the inverse of FKGL on a 0--100 scale
    (higher = more readable). More intuitive for non-specialist audiences and
    policy reporting.

  \item \textbf{Gunning Fog Index}: penalises words of three or more syllables,
    directly flagging the nominalisation problem identified in systemic functional
    linguistics (SFL) accounts of academic register
    \citep{halliday1993language}. A Gunning Fog score above 12 is considered
    difficult; above 17, obscure.

  \item \textbf{Lexical density}: the ratio of content words (nouns, verbs,
    adjectives, adverbs) to total words. In SFL, high lexical density characterises
    written, formal register; very high density in module descriptors may indicate
    information packing that reduces accessibility without increasing clarity
    \citep{halliday1989spoken}.

  \item \textbf{Average sentence length}: a direct measure of syntactic complexity;
    correlates with cognitive load independent of vocabulary.
\end{itemize}

Table~\ref{tab:readability} reports preliminary findings from the two institutions
for which the full pipeline has been run to date (MU and MTU; $n = 2{,}779$ and
$n = 6{,}110$ modules respectively with sufficient descriptor text).

\begin{table}[H]
\centering
\caption{Readability of module descriptors: MU (PU) vs MTU (TU). Higher FKGL and FOG
indicate less readable text; higher FRE indicates greater readability; LD = lexical
density (\%).}
\label{tab:readability}
\begin{tabular}{lrrrrr}
\toprule
Measure & \multicolumn{2}{c}{MU ($n=2{,}779$)} & \multicolumn{2}{c}{MTU ($n=6{,}110$)} \\
\cmidrule(lr){2-3}\cmidrule(lr){4-5}
 & Mean & Median & Mean & Median \\
\midrule
FKGL            & 18.74 & 17.53 & 20.14 & 17.82 \\
FRE             & 12.25 &  7.35 &  6.39 &  0.00 \\
Gunning Fog     & 21.35 & 20.12 & 22.75 & 20.44 \\
Lexical Density (\%) & 63.4 & 62.2 & 69.4 & 69.2 \\
Avg Sentence Length & 25.2 & 21.8 & 23.4 & 16.0 \\
\bottomrule
\end{tabular}
\end{table}

Both institutions produce module descriptors that are substantially harder to read
than is desirable for public-facing student documents. Mean FKGL values of 18.7 (MU)
and 20.1 (MTU) correspond to postgraduate reading-level requirements; the Gunning Fog
scores (21.4 and 22.7 respectively) far exceed the threshold of 17 typically considered
obscure. MTU descriptors show markedly higher lexical density than MU (69.4\% vs
63.4\%), consistent with the SFL prediction that professional/technical registers
sustain higher information packing \citep{halliday1989spoken} --- a finding that
warrants attention given the TU sector's stated commitment to student-centred learning.

We computed Pearson correlations between D1 verb tier (per module mean) and each
readability measure. The only substantive correlation was between D1 and lexical
density (MU: $r = 0.225$; MTU: $r = 0.191$), suggesting that modules specifying
more observable behaviours tend to use slightly denser, more content-rich prose ---
consistent with the hypothesis that precise LO writing requires richer noun-phrase
structures to specify conditions and criteria. Correlations between D1 and sentence-
level measures (FKGL, FOG, ASL) were negligible in both institutions ($|r| < 0.07$),
indicating that verb vagueness and syntactic complexity vary largely independently:
a module can be written in plain sentences yet specify unobservable outcomes, or use
complex prose to specify measurable ones.

\todo{Extend readability pipeline to all nine corpus institutions (DCU, Galway, UCC,
UCD, UL, ATU, SETU, TUD) once data collection is complete. Add cross-tab by NFQ level
and faculty/discipline. Consult Susan Gottl{\"o}ber re Humanities descriptor language
norms; Niamh Fortune re Social Sciences.}

% ── Beyond D1: Curriculum Design Signals ────────────────────
\subsubsection*{A note on the shadow curriculum}

The analyses in Sections~\ref{sec:udl}--\ref{sec:skills} share a methodological
caveat that is also a policy argument. Module descriptor text and published LOs
constitute the \textit{visible curriculum}: the documented commitment made to
students, accreditation bodies, and the public. They do not constitute the
\textit{enacted curriculum}. UDL practices, skills development activities, and
data literacy components that are present in teaching and assessment but absent
from documentation are invisible to this analysis --- and to every analysis
dependent on public-facing records.

This is not merely a measurement limitation. Good practice that goes undocumented
cannot be credited in national benchmarking, replicated by colleagues, or
communicated to prospective students making module choices. Where the analyses
below identify weak signal for UDL, skills, or data literacy, the most accurate
reading is: \textit{the documentation does not claim it}. Whether the practice
exists in the shadow curriculum is a question this dataset cannot answer --- but
it is one that institutions are well-placed to investigate and, where the evidence
supports it, to document visibly. The message is not deficit but opportunity:
if it is happening, tell everyone.

Where assessment component data is available (MTU, SETU, DkIT, NCI via Akari;
MU via CA/exam weighting), it provides a second window onto enacted practice
that the LO text alone cannot open. Sections below flag cases where assessment
data corroborates or contradicts the UDL, skills, or alignment signal in the
descriptor text.

\subsection{Universal Design for Learning: Proxy Analysis}
\label{sec:udl}

\note{This section reports preliminary UDL proxy findings for MU (2,139 modules).
Extension to the full corpus is planned. \textbf{[PLACEHOLDER: narrative to be
developed with Niamh Fortune and Susan Gottl{\"o}ber.]}}

Universal Design for Learning (UDL) \citep{cast2018udl} provides a complementary
lens to verb-level quality analysis. Where D1 asks what cognitive performance an LO
specifies, UDL principles ask whether the module's design creates the conditions for
diverse learners to achieve that performance. We operationalised three UDL principles
as proxies from module descriptor data at MU:

\begin{itemize}
  \item \textbf{Representation} (Principle I): diversity of teaching and learning
    modalities recorded in the timetabled activity (lectures, laboratory, tutorial,
    planned self-directed learning).
  \item \textbf{Action and Expression} (Principle II): diversity of assessment
    component types across the module.
  \item \textbf{Engagement} (Principle III): presence of engagement-signalling
    keywords in module descriptor text (e.g., \textit{collaborative},
    \textit{problem-based}, \textit{authentic}, \textit{reflective}).
\end{itemize}

Across 2,139 MU modules with complete data, preliminary findings indicate:

\begin{itemize}
  \item 40.3\% of modules report a single teaching and learning modality (lectures
    only), indicating low UDL Representation signal.
  \item Only 5.4\% of modules show strong signal across all three UDL principles.
  \item 17.9\% show no UDL signal (single T\&L modality, no engagement keywords,
    single assessment component).
  \item A weak positive correlation exists between UDL composite score and D1 mean
    ($r = 0.059$); engagement keyword presence shows the strongest relationship
    with verb observability ($r = 0.085$).
  \item Department-level variation is pronounced: modules in Critical Skills show the
    highest combined UDL and D1 profiles; modules in History show the lowest on
    both dimensions.
\end{itemize}

These findings are preliminary and subject to the limitations of keyword-based proxy
analysis. Descriptor text quality and completeness varies across departments and
individuals. \todo{Extend UDL analysis to full nine-institution corpus once descriptor
text is normalised across institutions. Consult with Niamh Fortune and Susan Gottl{\"o}ber re UDL in Irish HE context
(Froebel/Social Sciences and Humanities perspectives respectively). Consider whether Strand 4 UDL paper
material should be partially incorporated here or remain a separate publication.}

\subsection{HEA Data and Digital Skills: Module Tagging Case Study}
\label{sec:datadigital}

\note{\textbf{[PLACEHOLDER: data and narrative to be developed with Tim Thompson (VP for Students and Learning) and John G. Keating.]}}

In response to the HEA's \textit{Human Capital Initiative} mandate, Maynooth University
undertook a systematic tagging of modules contributing to Data and Digital Skills
competencies across its curriculum. This exercise combined human expert tagging with
AI-assisted classification, enabling an agreement analysis between human and
algorithmic judgements.

\todo{Develop this section once John shares HEA Data/Digital reports. Key content to
include: (a) methodology for human + AI tagging; (b) agreement analysis results;
(c) distribution of Data/Digital modules across faculties and NFQ levels; (d) connection
to LO quality findings --- does higher D1 verb tier correlate with stronger Data/Digital
signal? (e) relevance to national skills policy. This section connects Strands 3 and 5
of the broader research programme.}

\subsection{Graduate Skills Alignment: QQI Framework}
\label{sec:skills}

\note{\textbf{[PLACEHOLDER: analysis complete; narrative to be developed with
Máire Buckley, Experiential Learning Officer, MU.]}}

QQI's \textit{Graduate Attributes Framework} \citep{qqi2022graduates} identifies
a set of capabilities expected of Irish higher education graduates, including
critical thinking and problem solving, communication, collaboration, creativity
and innovation, research and inquiry, ethical responsibility, and self-direction.
These are not discipline-specific but constitute a cross-sectoral baseline against
which programme and module documentation can be examined.

We operationalise each attribute as a keyword and phrase set applied to both LO
text and module descriptor text across all fourteen institutions. Because
documentation is the unit of analysis (see shadow curriculum note above), the
measure captures whether institutions \textit{claim} these attributes in their
public module records, not whether they deliver them.

\begin{table}[H]
\centering
\caption{QQI Graduate Attribute keyword targets for corpus scanning.
Sources: LO text and module descriptor text; assessment component labels where
Akari data is available (MTU, SETU, DkIT, NCI, MU).}
\label{tab:skills-targets}
\begin{tabularx}{\textwidth}{lXl}
\toprule
\textbf{Attribute} & \textbf{LO / Descriptor targets} & \textbf{Assessment signal} \\
\midrule
Critical thinking & evaluate, critique, analyse, justify, argue, assess & essay, viva \\
Communication & present, report, communicate, write, articulate & presentation, report \\
Collaboration & collaborate, group work, team, peer review & group project \\
Creativity \& innovation & design, create, develop, innovate, generate & portfolio, project \\
Research \& inquiry & investigate, research, evidence-based, literature & dissertation, project \\
Ethical responsibility & ethical, responsible, sustainability, equity, justice & reflection \\
Self-direction & self-directed, independently, manage own learning & portfolio \\
\bottomrule
\end{tabularx}
\end{table}

\todo{Run scan-skills.mjs across all 14 corpus institutions. Produce: (a) per-institution
attribute coverage rates (proportion of modules with at least one keyword hit per
attribute); (b) sector-level comparison (PU vs TU vs IT vs CE vs PS); (c) cross-tab
with D1 tier (do high-D1 modules also show stronger skills signal?); (d) where Akari
assessment data is available, flag modules where assessment component type corroborates
the skills claim in the LO text. Consult Máire Buckley re QQI attribute definitions
and any MU-specific graduate attribute policy. Flag whether Data/Digital overlaps
with any QQI attribute (likely: research and inquiry, self-direction) to avoid
double-counting with Section~\ref{sec:datadigital}.}

\subsection{Alignment with National Policy Standards}
\label{sec:policyalignment}

\note{\textbf{[PLACEHOLDER: analysis complete; narrative and policy interpretation
to be developed with Lisa O'Regan, Centre for Teaching \& Learning, MU. See
\texttt{papers/policy-landscape-reference.md} for full policy language and
mapping table.]}}

The three-level policy framework described in Section~\ref{sec:policyframework}
establishes concrete, testable requirements for LO quality. Table~\ref{tab:policy-mapping}
maps these requirements to the D-dimensions of our rubric and summarises the
corpus-level compliance picture.

\begin{table}[H]
\centering
\caption{Mapping of policy requirements to rubric dimensions and corpus findings.
Sources: ESG 2015 Standard 1.2; QQI Validation Policy 2021; NFQ Grid of Level
Indicators; Dublin Descriptors (Cycle 1 and 2).}
\label{tab:policy-mapping}
\begin{tabularx}{\textwidth}{p{0.32\textwidth}p{0.22\textwidth}X}
\toprule
\textbf{Policy requirement} & \textbf{Rubric dimension} & \textbf{Corpus finding} \\
\midrule
\enquote{Specific and sufficiently detailed} (QQI) &
  D1 Verb observability &
  10.4\% of corpus LOs use Tier~0 (unobservable) verbs; range 2.6\%--23.9\% by institution \\
\addlinespace
\enquote{Specific} — one outcome per statement (implied) &
  D2 Behavioural singularity &
  \todo{D2 mean scores by institution — to insert} \\
\addlinespace
Student-centred LO format (ESG 1.2; Bologna) &
  D3 Student-centredness &
  TUD Vault (2,791 LOs) uses aims format; 736 student-centred LOs recovered from PDF catalogues across 6 programmes (TU835, TU874, TU078, TU079, TU248, TU5294); MIE 91.4\% of modules carry no student-centred LOs \\
\addlinespace
\enquote{Appropriate to specified NFQ level} (QQI) &
  D6 NFQ level calibration (LLM) &
  \todo{D6 results — to insert once LLM scoring complete} \\
\addlinespace
Assessment ensures MIPLOs are assessed (QQI) &
  Constructive alignment analysis &
  123 confirmed Type A modules across 5 institutions; MU: 49; MTU: 55; SETU: 7; DkIT: 10; NCI: 2 \\
\addlinespace
LOs present for all modules (QQI institutional QA) &
  LO publication rate &
  MIE: 8.6\% module LO rate; UCD: 31.2\% null; UL: 37.8\% null; TUD: aims-only \\
\bottomrule
\end{tabularx}
\end{table}

Several findings warrant particular attention from a policy compliance perspective.
First, the prevalence of Tier~0 verbs (\textit{understand}, \textit{know},
\textit{appreciate}) across all fourteen institutions indicates that despite
three decades of outcomes-based education policy, unobservable language remains
embedded in Irish HE module documentation at scale. Second, the MIE no-LO rate
(91.4\%) and TUD aims-format corpus represent structural rather than quality
issues: the problem is not poorly written LOs but the absence of the
student-centred outcome format that QQI and the Bologna process require. Third,
the 123 confirmed Type A modules identified across five institutions represent
documented cases in which QQI's assessment-alignment requirement is not met:
high-demand LOs are promised but examination-only assessment is applied,
producing a verifiable gap between the MIPLO commitment and the assessment
instrument designed to verify it.

\todo{Lisa O'Regan to develop: (a) contextualise findings against QQI's own
review/audit findings where publicly available; (b) consider whether the null
LO rates at UCD/UL represent institutional QA process failures or publication
system artefacts (and the policy implications of each); (c) institutional
response framing — recommendations for QA process improvement at programme
approval stage; (d) international comparator — how does Ireland compare with
EHEA peers on LO quality? Consider referencing EURASHE / EUA Learning and
Teaching Thematic Peer Group findings if available.}

% ════════════════════════════════════════════════════════════
\section{Discussion}
\label{sec:discussion}
% ════════════════════════════════════════════════════════════

\subsection{The PU--TU Gap}

The consistent D1 advantage of the TU sector over the public university cluster ---
$\Delta = 0.203$ (PU mean 1.572, TU mean 1.775), reproduced in direction across every
institution in both clusters --- merits careful interpretation.
It would be tempting to read this as evidence of a quality differential favouring TUs,
and in the specific dimension of verb observability this reading is defensible.
However, the interpretation should be contextualised.

TU programmes, rooted in applied and professional education, are designed around the
demonstration of competence in specific contexts. The programme approval and review
culture in the technological sector, historically regulated through QQI processes
with strong emphasis on measurable and vocational outcomes, may have produced a
curriculum documentation practice in which observable verbs are more consistently
selected. This is not necessarily evidence of deeper constructive alignment; it may
reflect a discipline-level norm about how outcomes are articulated.

Universities, by contrast, house a wider range of disciplines including humanities,
social sciences, and pure sciences in which the objectives of education are
genuinely more difficult to reduce to observable performance. A module in philosophy,
history, or theoretical mathematics may legitimately aim for transformative
\textit{understanding} that resists specification in the terms that Bloom's taxonomy
most easily accommodates. The persistent presence of \textit{understand} in university
LOs may therefore reflect both a quality problem (weak outcome design) and an
epistemological resistance (the genuine difficulty of specifying what excellent
philosophical or historical reasoning looks like in verb form).

The collaboration of Maynooth University colleagues in the authorship of this paper
reflects this tension directly. Faculty of Arts and Humanities and Faculty of Social
Sciences perspectives are essential for distinguishing the legitimate complexity of
humanistic learning objectives from the habitual use of vague language that could,
with appropriate design support, be made more precise without sacrificing
disciplinary integrity.

\subsection{Constructive Alignment: A Solvable Problem}

The finding that 113 modules across four institutions exhibit confirmed constructive
misalignment (MU: 49; MTU: 55; SETU: 7; NCI: 2) is notable not primarily because of
its scale, but because of what it implies about current quality assurance processes. These are not hidden or theoretical
misalignments: they are encoded in publicly accessible module descriptors and
institutional assessment records. A module approval process with routine access to
D1 verb scoring and assessment modality data would, in principle, detect and flag
these modules before students enrol.

The 101 MTU modules initially appearing misaligned but resolved by LO-index mapping
makes the complementary point: apparent misalignment is not always actual misalignment.
Surface-level analysis of exam percentages without reference to the LO-level mapping
of assessment components produces both false positives (modules flagged as misaligned
that are in fact well-designed) and false negatives (modules that appear acceptable
because their CA weight is high but whose CA is examination-like in structure, as the
SAQ finding suggests).

\subsection{Policy Implications}

Three policy implications follow from these findings.

\textbf{First}, module approval processes should incorporate routine LO quality
checking at the verb level. This does not require human expert review of 232,000
LOs; it requires a lightweight automated scoring step integrated into module
registration or approval systems. Our pipeline, which processes the full corpus in
under an hour on standard hardware, is deployable at institutional scale.

\textbf{Second}, assessment data recording standards should be revised to capture
component type, not merely percentage weight. The ambiguous category in MU's
analysis --- 24 modules that cannot be resolved without component-type data ---
is an artefact of data recording practice, not of the underlying educational reality.
The investment required to capture \textit{Project}, \textit{Short Answer Questions},
\textit{Presentation} alongside the percentage weight is minimal; the quality
assurance value is substantial.

\textbf{Third}, the TU Dublin finding invites a national conversation about the
publication of learning outcomes on public module catalogue pages. If module
descriptors are the primary instrument through which students make informed programme
choices, and through which the sector accounts for what it teaches, then the
widespread non-publication of student-centred learning outcomes is a transparency
issue that QQI's review processes should address directly.

\subsection{Limitations}

Several limitations should be noted. The D1--D3 dimensions capture syntactic and
lexical features of LO text; they do not assess the underlying quality of educational
design or the lived experience of alignment in the classroom. A module with perfectly
observable, student-centred, singularly specified LOs may still be poorly taught or
poorly assessed in practice; conversely, a module with Tier 0 verbs in its formal
descriptor may be taught with rigour and assessed authentically by a skilled practitioner.

The lexicon-based approach to D1 scoring is limited by lexicon coverage and by the
context-dependence of verb interpretation. \textit{Review}, for example, is assigned
Tier 1 in our lexicon, but in some disciplinary contexts it describes a complex
evaluative act more appropriately assigned to Tier 2. The v1.3 lexicon includes
a context-dependent verb register to flag such cases for manual review; a fully
automated treatment of context-dependence awaits LLM-based scoring (D4--D6, in progress).

The constructive alignment analysis at MU is limited by the weight-only recording of
continuous assessment. The 24 ambiguous modules cannot be resolved without
component-type data, and some proportion of the 49 confirmed Type A modules may
contain CA components not captured in the descriptor that do provide alignment rescue.
Similarly, the MTU analysis is limited to modules for which the Akari system records
LO-to-component mappings; a subset of modules had no such mapping and were excluded.

% ════════════════════════════════════════════════════════════
\section{Conclusion and Recommendations}
\label{sec:conclusion}
% ════════════════════════════════════════════════════════════

This study provides the first cross-institutional, corpus-scale characterisation of
learning outcome language quality in Irish higher education. Across 226,295 LOs from
nine institutions, we find a consistent gap in verb observability between university
and technological university sectors, with universities showing higher rates of
unobservable Tier 0 verbs and lower mean D1 scores. The gap is reproducible across
all institutions within each cluster and is unlikely to be artefactual.

The extension to constructive alignment reveals that verb quality is necessary but
not sufficient: 55 MTU modules (and a minimum of 49 MU modules) exhibit confirmed
misalignment between high-demand LO language and assessment modalities that cannot
elicit the specified performance. The methodological contrast between MU and MTU
demonstrates that the precision of this analysis is directly determined by the
granularity of assessment data recording practice.

We offer the following recommendations:

\begin{enumerate}
  \item \textbf{Integrate automated LO quality scoring into module approval workflows.}
    Lightweight rule-based scoring of the leading action verb is sufficient to flag
    Tier 0 usage for author review. This does not constrain academic freedom; it
    surfaces language choices for deliberate consideration.

  \item \textbf{Revise assessment data recording standards to capture component type.}
    The distinction between \textit{Project} and \textit{Short Answer Questions} as
    CA modalities is not a bureaucratic nicety; it is the difference between a
    resolvable and an unresolvable alignment question.

  \item \textbf{Establish a national expectation for LO publication on public module
    catalogue pages.} The TU Dublin finding indicates that the absence of student-centred
    LOs on public pages is not inevitable; it is an institutional choice. QQI programme
    review should include a check on public LO publication.

  \item \textbf{Extend the analysis to D4--D6.} Rule-based scoring captures syntax;
    LLM-based scoring of scope, assessability, and NFQ calibration will extend the
    analytical depth and move toward a composite quality score suitable for
    institutional dashboard reporting.

  \item \textbf{Involve academic development units in disciplinary interpretation.}
    The university--TU gap cannot be addressed by blanket verb-replacement advice.
    Humanities and social science disciplines require disciplinary-appropriate guidance
    on how to specify complex understanding in assessable terms without reducing
    learning to performance. Faculty-level academic development, informed by
    discipline-specific examples from the corpus, offers the most promising pathway.
\end{enumerate}

The data, verb lexicon, and scoring pipeline are available to institutional partners
and QQI on request. \todo{Add data availability / supplementary materials statement.}
\label{sec:data-availability}

% ════════════════════════════════════════════════════════════
\section*{Acknowledgements}
% ════════════════════════════════════════════════════════════

\todo{Add acknowledgements: MU ADTL office, participating institutions, any funding statement.}

% ════════════════════════════════════════════════════════════
\bibliographystyle{apalike}
\bibliography{lo-quality-refs}
% ════════════════════════════════════════════════════════════

% ════════════════════════════════════════════════════════════
\appendix
% ════════════════════════════════════════════════════════════

\section{LLM Evaluator Prompts}
\label{app:prompts}

This appendix reproduces the exact prompts used in the policy-grounded LLM
evaluation (\texttt{score-llm-v2.mjs}). Full scripts and all sixteen institutional
context blocks are available in the project repository
(see Section~\ref{sec:data-availability}).

% ── A.1 ───────────────────────────────────────────────────────────────────────
\subsection{System Prompt: Policy Layer and Evaluation Rubric}
\label{app:system-prompt}

The system prompt is fixed across all API calls and contains the complete policy
grounding and rubric (approximately 1,200 tokens; amortised across $N$ LOs per
call --- see Table~\ref{tab:batch}).

\begin{small}
\begin{verbatim}
You are an expert evaluator of learning outcomes (LOs) in Irish higher
education. Your evaluations are grounded in the following policy framework,
which defines what a well-formed LO must achieve in this jurisdiction.

## QQI Policies and Criteria for the Validation of Programmes (QP.17, 2017)
- Module LOs must be "explicitly specified" — vague or generalised
  statements are non-compliant.
- LOs must not be "vague, ambiguous or generalised."
- LOs must be "appropriate to award title and credit volumes."
- Minimum Intended Module Learning Outcomes (MIMLOs) are the binding
  specification.

## QQI Core Policies and Criteria (QP 24, 2024)
- LOs must be plainly written so that students can understand what they
  are expected to achieve.
- QQI's stated normative reference for LO writing guidance is Cedefop (2017).

## ECTS Users' Guide (European Commission, 2015)
- LOs describe what the learner will be able to know, understand, and do
  after completing a learning activity.
- LOs must describe demonstrable, observable achievement — not curriculum
  content coverage.
- Verbs such as "understand" and "know" are explicitly discouraged as
  unobservable.
- One ECTS credit corresponds to 25-30 hours of student workload — scope
  must be calibrated accordingly.

## Cedefop (2017) — Defining, Writing and Applying Learning Outcomes
   (QQI's normative reference)
- Each LO should specify a single observable behaviour using an active verb.
- Verbs should be drawn from a recognised taxonomy (e.g. Bloom's)
  appropriate to the cognitive level.
- LOs should be specific, measurable, and student-centred.
- Common errors: using vague verbs (know, understand, appreciate, be
  familiar with); writing aims or curriculum content as LOs; writing
  teacher-centred rather than learner-centred statements.

## NFQ Level Descriptors
NFQ Level 6: Knowledge — broad general knowledge of field; Skill —
  demonstrate a range of standard skills; Competence — act independently
  in familiar contexts, limited autonomy.
NFQ Level 7: Knowledge — broad knowledge with some theoretical depth;
  Skill — apply range of skills to known or new contexts; Competence —
  some autonomy in decision-making, take responsibility for outcomes.
NFQ Level 8 (Honours Bachelor): Knowledge — detailed theoretical and
  conceptual knowledge; Skill — demonstrate advanced skills including
  synthesis, evaluation, critical analysis; Competence — act with
  autonomy in novel complex contexts, take initiative.
NFQ Level 9 (Masters): Knowledge — specialised knowledge at forefront
  of field, including research; Skill — highly developed skills in
  research, analysis, and problem-solving in unpredictable contexts;
  Competence — act with substantial autonomy, manage complex professional
  activities.
NFQ Level 10 (Doctoral): Knowledge — original contribution to knowledge
  at forefront of field; Skill — expert research and advanced technical
  skills; Competence — exercise substantial authority, innovation, full
  autonomy in academic or professional activity.

## Evaluation Rubric

D1 — VERB OBSERVABILITY TIER (for IRR against NLP pipeline)
0 = Tier 0: Non-observable verb (know, understand, appreciate, be aware
    of, be familiar with, learn) — cannot be assessed as stated.
1 = Tier 1: Low-demand observable verb (identify, define, describe, list,
    name, state, recall) — assessable but cognitively undemanding.
2 = Tier 2: Mid-demand observable verb (apply, analyse, explain, compare,
    demonstrate, calculate, classify) — standard HE-level demand.
3 = Tier 3: High-demand observable verb (evaluate, critique, synthesise,
    design, create, construct, formulate, justify) — advanced demand.

D4 — SCOPE APPROPRIATENESS
Does the scope match the NFQ level and ECTS credit volume of the module?
0 = Scope is wrong: either a micro-task or a programme-level goal.
1 = Scope is plausible but imprecise.
2 = Scope is module-appropriate — a single assessable outcome at the
    right granularity for this credit weight.
3 = Scope is precisely calibrated — matches module content and credit
    weight exactly.

D5 — ASSESSABILITY SIGNAL
Can a specific, valid assessment be pictured for this outcome?
0 = No plausible assessment — describes a state, not a performance.
1 = An assessment can be imagined but requires significant reinterpretation.
2 = A clear assessment method is implied (exam question, lab task,
    report, presentation, portfolio).
3 = The LO directly signals a specific assessment approach, including
    context or artefact type.

D6 — NFQ LEVEL CALIBRATION
Is the cognitive demand appropriate for the stated NFQ level?
0 = Substantial mismatch — e.g. recall-only verb on NFQ 8/9.
1 = Mild mismatch — one level below or above expected cognitive demand.
2 = Consistent with NFQ level — the cognitive demand is appropriate.
3 = Precisely calibrated — verb and content complexity match both NFQ
    level and module context exactly.

Respond with ONLY valid JSON in exactly this schema:
{
  "d1_verify": { "score": <0-3>, "verb": "<extracted verb>",
                 "rationale": "<one sentence>" },
  "d4":        { "score": <0-3>, "rationale": "<one sentence>" },
  "d5":        { "score": <0-3>, "rationale": "<one sentence>" },
  "d6":        { "score": <0-3>, "rationale": "<one sentence>" }
}
\end{verbatim}
\end{small}

% ── A.2 ───────────────────────────────────────────────────────────────────────
\subsection{Representative User Prompt}
\label{app:user-prompt}

The user prompt is constructed per batch. The institutional context block is
selected by institution code from a lookup table of sixteen entries (one per
corpus institution); the module and LO fields are populated from the scored
corpus record. This example shows a single-LO call for clarity; in batch mode
the \texttt{LO~1}, \texttt{LO~2}, \ldots\ blocks are repeated within a single
user message and the model returns a JSON array of the same length.

\begin{small}
\begin{verbatim}
INSTITUTIONAL CONTEXT — MU:
Maynooth University: Action verb is mandatory — "the important word in
every learning outcome is the action verb." Staff required to write LOs
"appropriate to learners' level and stage." LOs should identify WHAT
the student does, not HOW. Maximum 8 LOs per module enforced by system.

LO 1 [AC205_LO1]:
Module: AC205 — MANAGEMENT & COST ACCOUNTING | NFQ Level 8 | 10 ECTS
Descriptor: The purpose of this course is to equip students with an
  explanation of management accounting as a basis for developing the
  tools and knowledge required to meet management information needs.
  Students will study the organisational content of management accounting;
  cost accumulation for inventory valuation and profit measurement;
  decision making; information for planning control and performance
  measurement, and ethics.
Text: "Describe and evaluate the management accountant's role,
  organizational environment, accounting information systems and
  possible ethical issues they face."

Evaluate D1_VERIFY, D4, D5, D6. Return JSON only.
\end{verbatim}
\end{small}

% ── A.3 ───────────────────────────────────────────────────────────────────────
\subsection{Output Schema and Example Response}
\label{app:output-schema}

The model returns one JSON object per LO (or a JSON array in batch mode). Scores
are clamped to $[0,3]$ and rounded to the nearest integer. Records that fail
JSON parsing or omit any required dimension are logged as errors and excluded
from analysis; they are retried automatically on the next run.

Example response for the LO shown in Appendix~\ref{app:user-prompt}:

\begin{small}
\begin{verbatim}
{
  "d1_verify": {
    "score": 1,
    "verb": "describe",
    "rationale": "The verb 'describe' is a low-demand observable verb,
                  making it assessable but not cognitively challenging."
  },
  "d4": {
    "score": 2,
    "rationale": "The scope is module-appropriate, addressing the role
                  of management accountants within the module context."
  },
  "d5": {
    "score": 2,
    "rationale": "A clear assessment is implied, such as an exam question
                  or report requiring description and evaluation of the
                  specified topics."
  },
  "d6": {
    "score": 1,
    "rationale": "Mild mismatch: the cognitive demand of 'describe' is
                  below what is expected for NFQ Level 8."
  }
}
\end{verbatim}
\end{small}

\noindent
The full set of sixteen institutional context blocks, all evaluation scripts,
and the complete scored output files are available at the project repository
(see Section~\ref{sec:data-availability}).

% ════════════════════════════════════════════════════════════
\section{Human Evaluation Rating Interface}
\label{app:rating-tool}

This appendix describes the browser-based rating tool used to collect human
evaluation data. The tool is publicly accessible at:

\begin{center}
\url{https://jkeatingmu.github.io/MU-finder/slo-study/}
\end{center}

\subsection{Architecture}

The application is a fully static single-page application (React/Vite), deployed
to GitHub Pages. It requires no authentication, no server, and no software
installation. Rating data are submitted directly from the browser to a Supabase
PostgreSQL instance via the Supabase REST API, using an insert-only Row Level
Security policy: anonymous users may insert rows but cannot read, update, or delete
any record. The researcher retrieves accumulated ratings via authenticated access to
the Supabase dashboard or by direct SQL query.

\subsection{Rating Workflow}

\begin{enumerate}
  \item The rater selects an institution from the sidebar. The full LO corpus for
  that institution loads from a pre-scored JSON file.
  \item The rater uses the search and filter controls to locate LOs of interest
  (by keyword, D1 tier, or NFQ level).
  \item Expanding an LO card reveals the full module context, the NLP measurement
  scores (D1--D3), and the LLM evaluation scores (D1\textsubscript{verify}, D4--D6),
  providing the rater with a reference point while maintaining independence.
  \item The rater clicks \textit{Rate this LO} to open the rating modal.
  \item The modal presents the LO text at the top. A collapsible \textit{Show
  module descriptor} toggle is available below the LO text for raters who need
  the broader module context (particularly relevant for D4--D6).
  \item The modal solicits a score and optional rationale for all six dimensions
  D1--D6 in sequence. Scores are entered via a labelled range slider (0--3 for
  D1/D4/D5/D6; 1--3 for D2/D3). The rubric descriptor for the selected score
  updates dynamically below each slider, colour-coded red/orange/green/blue for
  scores 0/1/2/3. Rationale text boxes include a prompt not to enter personal
  information. A final \textit{Confidence in these ratings} slider (1--5) records
  methodological metadata about rater certainty.
  \item On submission, the record is inserted into the \texttt{lo\_ratings} table
  with a server-side timestamp. The rater's browser records the \texttt{lo\_id} in
  \texttt{localStorage}; rated cards display a \textit{Rated} indicator on
  subsequent visits.
\end{enumerate}

\subsection{Data Schema}

The \texttt{lo\_ratings} table schema is as follows:

\begin{center}
\begin{tabular}{lll}
\toprule
Column & Type & Description \\
\midrule
\texttt{id} & UUID & Primary key (auto-generated) \\
\texttt{created\_at} & timestamptz & Submission timestamp (server-side) \\
\texttt{lo\_id} & text & Unique LO identifier (corpus key) \\
\texttt{institution} & text & Institution code (e.g.\ MU, DCU) \\
\texttt{module\_code} & text & Module code (e.g.\ AC205) \\
\texttt{lo\_text} & text & Full LO text (for audit) \\
\texttt{nfq\_level} & integer & NFQ level of the module \\
\texttt{rater} & text & \textit{anonymous} (pending ethics approval) \\
\texttt{d1} & integer [0--3] & D1 verb observability score \\
\texttt{d1\_rationale} & text & Optional rationale \\
\texttt{d2} & integer [1--3] & D2 behavioural singularity score \\
\texttt{d2\_rationale} & text & Optional rationale \\
\texttt{d3} & integer [1--3] & D3 student-centredness score \\
\texttt{d3\_rationale} & text & Optional rationale \\
\texttt{d4} & integer [0--3] & D4 scope score \\
\texttt{d4\_rationale} & text & Optional rationale \\
\texttt{d5} & integer [0--3] & D5 assessability score \\
\texttt{d5\_rationale} & text & Optional rationale \\
\texttt{d6} & integer [0--3] & D6 NFQ calibration score \\
\texttt{d6\_rationale} & text & Optional rationale \\
\texttt{confidence} & integer [1--5] & Rater's self-assessed confidence \\
\bottomrule
\end{tabular}
\end{center}

\noindent
Rater identity is currently recorded as \textit{anonymous} pending institutional
ethics approval for identified data collection. The schema accommodates a named
\texttt{rater} field; authenticated rating (e.g.\ via institutional SSO or a
shared session token) is a planned extension that will be activated once ethical
clearance is in place.

\end{document}
